\documentclass[11pt,a4paper]{article}
\usepackage[utf8]{inputenc}
\usepackage[T1]{fontenc}
\usepackage[catalan]{babel}
\usepackage{lmodern}
\linespread{1.22}
\usepackage[a4paper,top=2.4cm,bottom=2.2cm,left=1.5cm,right=1.5cm,headheight=1.1cm,headsep=0.5cm,footskip=1cm]{geometry}
\usepackage{amsmath,amssymb,cancel}
\usepackage{icomma}
\usepackage{graphicx}
\usepackage[table]{xcolor}
\usepackage{tikz}\usetikzlibrary{calc}
\usepackage{pgfplots}\pgfplotsset{compat=1.17}
\usepackage{fontawesome5}
\usepackage{ccicons}
\usepackage{fancyhdr}
\usepackage{enumitem}
\usepackage{booktabs}
\usepackage{tabularx}
\usepackage{needspace}
\usepackage[most]{tcolorbox}
\usepackage{titlesec}
\usepackage[hidelinks]{hyperref}
\hypersetup{pdftitle={Apunts de teoria. Unitat 5: Límits i continuïtat. 2n de Batxillerat MII},pdfsubject={Matemàtiques II},pdfkeywords={límits, indeterminacions, límits trigonomètrics, infinitèsims equivalents, paràmetres, continuïtat, funcions a trossos, teorema de Bolzano, asímptotes}}

% colors de la web (FentMates, bloc d'Anàlisi: verd; es manté el nom fmgreen dels apunts de les altres unitats)
\definecolor{fmnavy}{HTML}{1B2748}
\definecolor{fmgreen}{HTML}{6A8A33}
\definecolor{fmgreenl}{HTML}{7FA040}
\definecolor{fmyellow}{HTML}{FFC94D}
\definecolor{fmgold}{HTML}{E9A919}
\definecolor{fmaccent}{HTML}{FF7A50}
\definecolor{fmteal}{HTML}{2BB3A3}
\definecolor{fmdef}{HTML}{1D4F73}\definecolor{fmdefbg}{HTML}{EEF4F9}
\definecolor{fmprop}{HTML}{3B4252}\definecolor{fmpropbg}{HTML}{F3F4F7}
\definecolor{fmex}{HTML}{A3203A}\definecolor{fmexbg}{HTML}{FBF0F2}
\definecolor{fmmuted}{HTML}{5D6685}\definecolor{fmball0}{HTML}{E2483D}\definecolor{fmball1}{HTML}{2F6FD6}
\definecolor{fmsurface}{HTML}{EEF1F8}

\newcommand{\lt}{<}\newcommand{\gt}{>}
% fórmula en línia que no es parteix mai; si no cap en una línia, es posa centrada
\newcommand{\fmm}[1]{\setbox0=\hbox{$#1$}\ifdim\wd0>\linewidth\fmdisp{#1}\else\mbox{\box0}\fi}
% fórmula centrada; si és més ampla que la línia, s'encongeix fins que hi cap
\newcommand{\fmdisp}[1]{\[\setbox2=\hbox{$\displaystyle #1$}\ifdim\wd2>\linewidth\resizebox{\linewidth}{!}{\box2}\else\box2\fi\]}
\DeclareUnicodeCharacter{2212}{-}
\DeclareUnicodeCharacter{2225}{\ensuremath{\parallel}}\DeclareUnicodeCharacter{22A5}{\ensuremath{\perp}}\DeclareUnicodeCharacter{2226}{\ensuremath{\nparallel}}
\DeclareUnicodeCharacter{207A}{\ensuremath{^{+}}}\DeclareUnicodeCharacter{207B}{\ensuremath{^{-}}}
\setlength{\parindent}{0pt}
\setlength{\parskip}{12pt}
\setlength{\lineskip}{5pt}\setlength{\lineskiplimit}{5pt}  % més aire al voltant de les fraccions en línia
\setlength{\emergencystretch}{3em}
\setlist{itemsep=4pt,topsep=4pt}
\renewcommand{\thesection}{5.\arabic{section}}
\titleformat{\section}{\color{fmnavy}\Large\bfseries}{\colorbox{fmdef}{\color{white}\strut\ \thesection\ }}{0.6em}{}[{\color{fmdef}\titlerule[1.2pt]}]
\titlespacing*{\section}{0pt}{16pt}{8pt}
\titleformat*{\subsection}{\color{fmgreen}\large\bfseries}
\allowdisplaybreaks

% dibuix de la targeta de la unitat a la web (illustrations.js), en blanc
\newcommand{\fmcard}{\begin{tikzpicture}[x=0.062cm,y=-0.062cm,white,line cap=round]
  \begin{scope}[transparency group,opacity=0.15]
  \draw[line width=3pt] (14,104)--(226,104) (24,114)--(24,8);
  \draw[line width=4pt,dash pattern=on 10pt off 10pt] (24,36)--(226,36);
  \draw[line width=4.5pt] (32,100) .. controls (70,98) and (92,80) .. (112,64);
  \draw[line width=4.5pt] (122,58) .. controls (150,44) and (180,40) .. (222,38);
  \draw[line width=4.5pt] (117,61) circle (5.5);
  \node[anchor=base west,inner sep=0pt,scale=2.4,font=\large\itshape] at (138,84) {lím $f(x)=L$};
  \end{scope}
\end{tikzpicture}}
\newcommand{\fmbrand}{\raisebox{-0.25\height}{\includegraphics[height=0.5cm]{logo-fentmates.pdf}}}

% capçalera i peu
\pagestyle{fancy}\fancyhf{}
\fancyhead[L]{\small\fmbrand}
\fancyhead[R]{\small\color{fmmuted}2n Batxillerat MII · Unitat 5: Límits i continuïtat}
\fancyfoot[L]{\scriptsize\color{fmmuted}\href{https://creativecommons.org/licenses/by-nc-sa/4.0/deed.ca}{\ccbyncsa}\ \ FentMates · fentmates.org · CC BY-NC-SA 4.0}
\fancyfoot[R]{\small\color{fmmuted}\thepage}
\renewcommand{\headrulewidth}{0.4pt}\renewcommand{\headrule}{\color{fmmuted!40}\hrule height \headrulewidth}
\fancypagestyle{primera}{\fancyhf{}\renewcommand{\headrulewidth}{0pt}
  \fancyfoot[L]{\scriptsize\color{fmmuted}\href{https://creativecommons.org/licenses/by-nc-sa/4.0/deed.ca}{\ccbyncsa}\ \ FentMates · fentmates.org · CC BY-NC-SA 4.0}
  \fancyfoot[R]{\small\color{fmmuted}\thepage}}

% caixes (mateixos colors que la web)
% quan una caixa es parteix entre dues pàgines, la vora del tall és discontínua
\tcbset{fmbase/.style={enhanced,breakable,arc=3mm,boxrule=1pt,left=3.5mm,right=3.5mm,top=2mm,bottom=2.5mm,before skip=10pt,after skip=10pt,fonttitle=\bfseries,
  frame style={fill=tcbcolframe},
  overlay first={\fill[tcbcolback] ([xshift=1pt]frame.south west) rectangle ([xshift=-1pt,yshift=1.2pt]frame.south east);
    \draw[tcbcolframe,line width=1pt,dash pattern=on 5pt off 3pt] (frame.south west)--(frame.south east);},
  overlay middle={\fill[tcbcolback] ([xshift=1pt]frame.south west) rectangle ([xshift=-1pt,yshift=1.2pt]frame.south east);
    \fill[tcbcolback] ([xshift=1pt]frame.north west) rectangle ([xshift=-1pt,yshift=-1.2pt]frame.north east);
    \draw[tcbcolframe,line width=1pt,dash pattern=on 5pt off 3pt] (frame.south west)--(frame.south east) (frame.north west)--(frame.north east);},
  overlay last={\fill[tcbcolback] ([xshift=1pt]frame.north west) rectangle ([xshift=-1pt,yshift=-1.2pt]frame.north east);
    \draw[tcbcolframe,line width=1pt,dash pattern=on 5pt off 3pt] (frame.north west)--(frame.north east);}}}
\newtcolorbox{fmdef}[1]{fmbase,colback=fmdefbg,colframe=fmdef,coltitle=white,title={#1}}
\newtcolorbox{fmprop}[1]{fmbase,colback=fmpropbg,colframe=fmprop,coltitle=white,title={#1},halign title=center}
\newtcolorbox{fmex}[1]{fmbase,before upper=\raggedright,colback=fmexbg,colframe=fmex,colbacktitle=fmexbg,coltitle=fmex,title={\faPencil*\ #1},
  titlerule=0.4pt,titlerule style={fmex!40,dashed}}
\newtcolorbox{fmwarn}[1]{fmbase,colback=fmsf!6,colframe=fmsf,coltitle=white,title={#1}}
\newtcolorbox{fmexfill}[1]{fmbase,breakable=false,height fill,colback=fmexbg,colframe=fmex,colbacktitle=fmexbg,coltitle=fmex,title={\faPencil*\ #1},
  titlerule=0.4pt,titlerule style={fmex!40,dashed}}
\newtcolorbox{fmnote}{enhanced,breakable,colback=white,colframe=white,borderline west={3pt}{0pt}{fmaccent},arc=0pt,left=3mm,right=2mm,top=1mm,bottom=1mm,before skip=6pt,after skip=6pt,before upper=\raggedright}
\newcommand{\fmkey}[1]{\colorbox{fmyellow!45}{\strut #1}}
\newcommand{\fmuline}[1]{\textbf{\color{fmaccent}#1}}
\definecolor{fmok}{HTML}{15803D}\definecolor{fmev}{HTML}{2563EB}\definecolor{fmsf}{HTML}{C2410C}\definecolor{fmsi}{HTML}{7C3AED}\definecolor{fmas}{HTML}{A3203A}
\newcommand{\fmbadge}[2]{\tcbox[on line,colback=#1,colframe=#1,arc=2mm,boxsep=0pt,left=4pt,right=4pt,top=1.5pt,bottom=1.5pt]{\color{white}\footnotesize\bfseries #2}}
\newcommand{\fmchip}[1]{\tcbox[on line,colback=fmsf!10,colframe=fmsf!50,arc=2mm,boxsep=1pt,left=4pt,right=4pt]{#1}}
\newcommand{\fmlevel}[2]{\tcbox[on line,colback=white,colframe=fmmuted!40,arc=2mm,boxsep=2pt,halign=center]{\begin{tabular}{c}{\scriptsize\color{fmmuted}\MakeUppercase{#1}}\\#2\end{tabular}}}
% gràfiques
\pgfplotsset{fmaxis/.style={axis lines=middle,grid=major,grid style={fmmuted!15},axis line style={fmmuted,-latex},
  tick label style={font=\scriptsize,color=fmmuted},clip=true,unbounded coords=jump,restrict y to domain=-200:200,
  every axis plot/.append style={line join=round}},
  fmcurve/.style={fmaccent,line width=1.6pt,smooth}}
\tikzset{fmasy/.style={fmmuted,dashed,line width=0.9pt},fmguide/.style={fmaccent,dotted,line width=0.8pt},
  fmdot/.style={circle,fill=fmaccent,inner sep=0pt,minimum size=5pt},fmopen/.style={circle,draw=fmaccent,fill=white,line width=1pt,inner sep=0pt,minimum size=5pt}}
\newcommand{\fmpau}[1]{\tcbox[on line,colback=fmnavy,colframe=fmnavy,arc=2mm,boxsep=0pt,left=4pt,right=4pt,top=1.5pt,bottom=1.5pt]{\color{white}\scriptsize\bfseries\MakeUppercase{#1}}\ }

\begin{document}
\thispagestyle{primera}
\begin{tikzpicture}[remember picture,overlay]
  \fill[fmgreen] (current page.north west) rectangle ([yshift=-10.1cm]current page.north east);
  \begin{scope}
    \clip (current page.north west) rectangle ([yshift=-10.1cm]current page.north east);
    \node[anchor=north west,inner sep=0pt] at ([xshift=6.6cm,yshift=-1.3cm]current page.north west) {\fmcard};
  \end{scope}
  \node[anchor=north west,inner sep=0pt] at ([xshift=1.5cm,yshift=-1.1cm]current page.north west) {\tcbox[colback=white,colframe=white,arc=2.5mm,boxsep=0pt,left=7pt,right=7pt,top=5pt,bottom=5pt]{\includegraphics[height=0.75cm]{logo-fentmates.pdf}}};
\end{tikzpicture}
\vspace*{0.55cm}
\tcbox[colback=white!15!fmgreen,colframe=white!15!fmgreen,arc=2.5mm,boxsep=0pt,left=6pt,right=6pt,top=3pt,bottom=3pt]{\color{white}\small\bfseries Apunts de teoria · 2n Batxillerat · Matemàtiques II}

\vspace{4pt}
{\raggedright\color{white}\fontsize{25}{30}\selectfont\bfseries{\color{fmyellow}Unitat 5} · Límits i continuïtat\par}
\vspace{6pt}
\begin{minipage}{\linewidth}\color{white!88!fmgreen}\small Teoria i exemples resolts pas a pas: límits a l'infinit i en un nombre, totes les indeterminacions (\fmm{\frac{\infty}{\infty}}, \fmm{\infty-\infty}, \fmm{0\cdot\infty}, \fmm{1^\infty}, \fmm{\frac00}), límits trigonomètrics i infinitèsims equivalents, límits que depenen d'un paràmetre, continuïtat de funcions a trossos amb trossos polinòmics, racionals, irracionals, exponencials, logarítmics i trigonomètrics, teorema de Bolzano i asímptotes de qualsevol funció.\end{minipage}\par
% l'índex comença sempre a la mateixa altura, per davall de la franja (com en les unitats 5-10)
\vspace*{\dimexpr 8.2cm-\pagetotal\relax}

\begin{tcolorbox}[colback=fmsurface,colframe=fmsurface,arc=3mm,left=4mm,right=4mm,top=2mm,bottom=2mm]
\small\textbf{Índex}\par\smallskip
\begin{minipage}[t]{0.49\linewidth}\small \mbox{\hyperlink{sec1}{\textbf{\color{fmgreen}5.1}\ \ Límits a l'infinit: lectura gràfica}}\\[3pt]
\mbox{\hyperlink{sec2}{\textbf{\color{fmgreen}5.2}\ \ Càlcul analític a l'infinit}}\\[3pt]
\mbox{\hyperlink{sec3}{\textbf{\color{fmgreen}5.3}\ \ Comparació d'infinits}}\\[3pt]
\mbox{\hyperlink{sec4}{\textbf{\color{fmgreen}5.4}\ \ Indeterminacions a l'infinit}}\\[3pt]
\mbox{\hyperlink{sec5}{\textbf{\color{fmgreen}5.5}\ \ Límits en un nombre}}\\[3pt]
\mbox{\hyperlink{sec6}{\textbf{\color{fmgreen}5.6}\ \ Indeterminacions en un nombre}}\\[3pt]
\mbox{\hyperlink{sec7}{\textbf{\color{fmgreen}5.7}\ \ Límits trigonomètrics i infinitèsims}}\\[3pt]\end{minipage}\hfill\begin{minipage}[t]{0.49\linewidth}\small \mbox{\hyperlink{sec8}{\textbf{\color{fmgreen}5.8}\ \ Límits amb paràmetres}}\\[3pt]
\mbox{\hyperlink{sec9}{\textbf{\color{fmgreen}5.9}\ \ Continuïtat en un punt}}\\[3pt]
\mbox{\hyperlink{sec10}{\textbf{\color{fmgreen}5.10}\ \ Continuïtat de funcions a trossos}}\\[3pt]
\mbox{\hyperlink{sec11}{\textbf{\color{fmgreen}5.11}\ \ Funcions a trossos amb paràmetres}}\\[3pt]
\mbox{\hyperlink{sec12}{\textbf{\color{fmgreen}5.12}\ \ Teorema de Bolzano}}\\[3pt]
\mbox{\hyperlink{sec13}{\textbf{\color{fmgreen}5.13}\ \ Asímptotes}}\\[3pt]\end{minipage}
\end{tcolorbox}

\Needspace{4\baselineskip}
\section{Límits a l'infinit: càlcul geomètric}\hypertarget{sec1}{}
\Needspace{3\baselineskip}
\begin{fmdef}{Definició: límit quan x tendeix a infinit}
Calcular el límit d'una funció \fmm{f(x)} quan \fmm{x} tendeix a l'infinit és conéixer el comportament d'aquesta funció quan \fmm{x} es fa \textbf{tan gran com vulguem} (\fmm{x\to+\infty}) o \textbf{tan negativa com vulguem} (\fmm{x\to-\infty}). Podem trobar tres possibilitats:

\begin{itemize}
\item \fmm{\displaystyle\lim_{x\to+\infty}f(x)=L}: els valors de \fmm{f(x)} s'acosten tant com vulguem al nombre \fmm{L}. Gràficament, la corba s'acosta a la recta horitzontal \fmm{y=L}, que és una \textbf{asímptota horitzontal}.
\item \fmm{\displaystyle\lim_{x\to+\infty}f(x)=+\infty} (o \fmm{-\infty}): els valors de \fmm{f(x)} creixen (o decreixen) sense fita: la corba puja (o baixa) indefinidament.
\item \textbf{No existeix} el límit: els valors de \fmm{f(x)} no s'estabilitzen al voltant d'un nombre ni creixen o decreixen cap a l'infinit. També pot passar que la funció no estiga definida en algun dels infinits, és a dir, que el seu domini comence o acabe en un nombre determinat.
\end{itemize}

El mateix val quan \fmm{x\to-\infty}, mirant la part \textbf{esquerra} de la gràfica.
\end{fmdef}

\Needspace{7\baselineskip}
\begin{fmexfill}{Exemples: llegim els límits a la gràfica}
\begin{minipage}[t][\tcbtextheight][s]{\linewidth}
\noindent \begin{minipage}[t]{0.48\linewidth}\centering $f(x)=\dfrac{2x^2}{x^2+1}$\par\smallskip \begin{tikzpicture}\begin{axis}[fmaxis,width=6.4cm,height=4.6cm,xmin=-6,xmax=6,ymin=-1,ymax=3.2]\addplot[fmcurve,domain=-6:6,samples=160]{2*x^2/(x^2+1)};\draw[fmasy] ({rel axis cs:0,0}|-{axis cs:0,2}) -- ({rel axis cs:1,0}|-{axis cs:0,2});\end{axis}\end{tikzpicture}\par\smallskip $\displaystyle\lim_{x\to+\infty}f(x)=2,\qquad\lim_{x\to-\infty}f(x)=2$ \\ \vspace*{0.2cm} Asímptota horitzontal $y=2$ pels dos costats.\end{minipage}\hfill\begin{minipage}[t]{0.48\linewidth}\centering $g(x)=2^x-1$\par\smallskip \begin{tikzpicture}\begin{axis}[fmaxis,width=6.4cm,height=4.6cm,xmin=-6,xmax=4,ymin=-2,ymax=6]\addplot[fmcurve,domain=-6:3,samples=160]{2^x-1};\draw[fmasy] ({rel axis cs:0,0}|-{axis cs:0,-1}) -- ({rel axis cs:1,0}|-{axis cs:0,-1});\end{axis}\end{tikzpicture}\par\smallskip $\displaystyle\lim_{x\to-\infty}g(x)=-1,\qquad\lim_{x\to+\infty}g(x)=+\infty$ \\ \vspace*{0.2cm} Asímptota horitzontal $y=-1$ només per l’esquerra.\end{minipage}\par
\vfill
\noindent \begin{minipage}[t]{0.48\linewidth}\centering $h(x)=\dfrac{x^3}{10}-x$\par\smallskip \begin{tikzpicture}\begin{axis}[fmaxis,width=6.4cm,height=4.6cm,xmin=-6,xmax=6,ymin=-6,ymax=6]\addplot[fmcurve,domain=-6:6,samples=160]{x^3/10-x};\end{axis}\end{tikzpicture}\par\smallskip $\displaystyle\lim_{x\to-\infty}h(x)=-\infty,\qquad\lim_{x\to+\infty}h(x)=+\infty$\\ \vspace*{0.2cm} No té asímptota horitzontal.\end{minipage}\hfill\begin{minipage}[t]{0.48\linewidth}\centering $k(x)=\sin x$\par\smallskip \begin{tikzpicture}\begin{axis}[fmaxis,width=6.4cm,height=4.6cm,xmin=-10,xmax=10,ymin=-2,ymax=2]\addplot[fmcurve,domain=-10:10,samples=300]{sin(deg(x))};\end{axis}\end{tikzpicture}\par\smallskip $\displaystyle\nexists\lim_{x\to-\infty}\sin x,\qquad\nexists\lim_{x\to+\infty}\sin x$\\ \vspace*{0.2cm} Osci\lgem{}a entre $-1$ i $1$ sense parar.\end{minipage}\par
\vfill
\noindent \begin{minipage}[t]{0.48\linewidth}\centering $m(x)=\sqrt{x-2}$\par\smallskip \begin{tikzpicture}\begin{axis}[fmaxis,width=6.4cm,height=4.6cm,xmin=-4,xmax=10,ymin=-1,ymax=4]\addplot[fmcurve,domain=2:10,samples=300]{sqrt(max(x-2,0))};\node[fmdot] at (axis cs:2,0) {};\end{axis}\end{tikzpicture}\par\smallskip $\displaystyle\nexists\lim_{x\to-\infty}m(x),\qquad\lim_{x\to+\infty}m(x)=+\infty$\\ \vspace*{0.2cm} Domini $[2,+\infty)$: no està definida per a $x<2$.\end{minipage}\hfill\begin{minipage}[t]{0.48\linewidth}\centering $n(x)=\sqrt{9-x^2}$\par\smallskip \begin{tikzpicture}\begin{axis}[fmaxis,width=6.4cm,height=4.6cm,xmin=-5,xmax=5,ymin=-1,ymax=4]\addplot[fmcurve,domain=-3:3,samples=300]{sqrt(max(9-x^2,0))};\node[fmdot] at (axis cs:-3,0) {};\node[fmdot] at (axis cs:3,0) {};\end{axis}\end{tikzpicture}\par\smallskip $\displaystyle\nexists\lim_{x\to-\infty}n(x),\qquad\nexists\lim_{x\to+\infty}n(x)$\\ \vspace*{0.2cm} Domini $[-3,3]$: no arriba a cap dels dos infinits.\end{minipage}\par
\vfill

      
\begin{fmnote}\fmuline{Truc}: per a \fmm{x\to+\infty} mira \textbf{el final dret} de la gràfica; per a \fmm{x\to-\infty}, \textbf{el final esquerre}. Una taula de valors amb \fmm{x=10,\ 100,\ 1000,\dots} també t'ho diu.\end{fmnote}
\end{minipage}
\end{fmexfill}

\Needspace{18\baselineskip}
\section{Càlcul analític de límits a l'infinit}\hypertarget{sec2}{}

\Needspace{7\baselineskip}
\begin{fmprop}{Operacions amb infinit}
Per calcular un límit, \textbf{substituïm} i operem amb aquestes regles (\fmm{k} és un nombre real):

\begin{center}\small\begin{tabularx}{\linewidth}{>{\centering\arraybackslash}X>{\centering\arraybackslash}X>{\centering\arraybackslash}X}\toprule
\textbf{Sumes} & \textbf{Productes} & \textbf{Quocients}\\\midrule
\fmm{k\pm\infty=\pm\infty} & \fmm{k\cdot(\pm\infty)=\pm\infty\ \ (k\gt0)} & \fmm{\dfrac{k}{\pm\infty}=0}\\[6pt]
\fmm{(+\infty)+(+\infty)=+\infty} & \fmm{k\cdot(\pm\infty)=\mp\infty\ \ (k\lt0)} & \fmm{\dfrac{\pm\infty}{k}=\pm\infty\ \ (k\gt0)}\\[6pt]
\fmm{(-\infty)+(-\infty)=-\infty} & \fmm{(\pm\infty)\cdot(\pm\infty)=\pm\infty} & \fmm{\dfrac{k}{0}=\pm\infty\ \ (k\neq0)} \\[6pt]
&(regla dels signes)&(cal estudiar el signe)\\[6pt]
\bottomrule\end{tabularx}\end{center}

\par\noindent \begin{minipage}[t]{0.48\linewidth}\vspace{0pt}\raggedright \textbf{Exponencials}

\fmdisp{k^{+\infty}=\begin{cases}+\infty&\text{si } k\gt 1\\0&\text{si } 0\leq k\lt 1\end{cases}}

\fmdisp{k^{-\infty}=\begin{cases}0&\text{si } k\gt 1\\+\infty&\text{si } 0\lt k\lt 1\end{cases}}
\end{minipage}\hfill\begin{minipage}[t]{0.48\linewidth}\vspace{0pt}\raggedright \textbf{Potències}

\fmdisp{(+\infty)^k=\begin{cases}+\infty&\text{si } k\gt 0\\0&\text{si } k\lt 0\end{cases}}

\fmdisp{(+\infty)^{+\infty}=+\infty,\ \ (+\infty)^{-\infty}=0}
\end{minipage}\par\medskip

Hi ha resultats que \textbf{no es poden saber} d'entrada: són les \textbf{indeterminacions}, que estudiarem a l'apartat 5.4:

\begin{center}\fmchip{\fmm{\dfrac{\infty}{\infty}}}\ \fmchip{\fmm{\infty-\infty}}\ \fmchip{\fmm{\dfrac00}}\ \fmchip{\fmm{0\cdot\infty}}\ \fmchip{\fmm{1^{\infty}}}\ \fmchip{\fmm{\infty^0}}\ \fmchip{\fmm{0^0}}\end{center}
\end{fmprop}

\Needspace{7\baselineskip}
\begin{fmdef}{Potències, monomis i polinomis}
\begin{itemize}
\item \textbf{Potències} \fmm{x^n} amb \fmm{n\gt0}: quan \fmm{x\to+\infty} sempre donen \fmm{+\infty}. Quan \fmm{x\to-\infty}, depén de la paritat:

\fmdisp{\lim_{x\to-\infty}x^n=\begin{cases}+\infty&\text{si } n \text{ parell}\\-\infty&\text{si } n\text{ senar}\end{cases}\qquad\qquad\lim_{x\to\pm\infty}\frac{1}{x^n}=0}

\item \textbf{Monomis} \fmm{a\,x^n}: el límit és \fmm{\pm\infty}; el signe el decideixen el signe de \fmm{a} i el de \fmm{x^n}.
\item \textbf{Polinomis}: quan \fmm{x\to\pm\infty}, el límit d'un polinomi és el límit del seu \fmkey{terme de major grau}; els altres termes són «petits» al costat d'ell.

\fmdisp{\lim_{x\to\pm\infty}\left(a_nx^n+a_{n-1}x^{n-1}+\dots+a_0\right)=\lim_{x\to\pm\infty}a_nx^n}

\end{itemize}

\begin{fmnote}\fmuline{Compte}: \fmm{(-\infty)^3} és \fmm{-\infty}, però \fmm{(-\infty)^4} és \fmm{+\infty}. Si dubtes, canvia \fmm{x} per \fmm{-x}: \fmm{\displaystyle\lim_{x\to-\infty}f(x)=\lim_{x\to+\infty}f(-x)}.\end{fmnote}
\end{fmdef}

\Needspace{7\baselineskip}
\begin{fmex}{Exemples}
\begin{enumerate}[label=\textbf{\arabic*.},leftmargin=*]
\item \fmm{\displaystyle\lim_{x\to-\infty}2x^3=2\cdot(-\infty)^3=2\cdot(-\infty)=-\infty}  i  \fmm{\displaystyle\lim_{x\to-\infty}(-5x^3)=(-5)\cdot(-\infty)=+\infty}.
\item \fmm{\displaystyle\lim_{x\to+\infty}\left(x^3-100x^2-7\right)=\lim_{x\to+\infty}x^3=+\infty}. Encara que \fmm{100x^2} sembla gran, \fmm{x^3} guanya.
\item \fmm{\displaystyle\lim_{x\to-\infty}\left(-2x^4+x^3+1\right)=\lim_{x\to-\infty}(-2x^4)=-2\cdot(+\infty)=-\infty}.
\item \fmm{\displaystyle\lim_{x\to+\infty}\left(\tfrac12\right)^x=\left(\tfrac12\right)^{+\infty}=0}, \fmm{\ \ \displaystyle\lim_{x\to-\infty}3^x=3^{-\infty}=0}, \fmm{\ \ \displaystyle\lim_{x\to+\infty}\sqrt{x}=\sqrt{+\infty}=+\infty}.
\end{enumerate}
\end{fmex}

\Needspace{7\baselineskip}
\begin{fmdef}{Quocients de polinomis: comparem els graus}
Si \fmm{P(x)} té grau \fmm{n} i coeficient principal \fmm{a_n}, i \fmm{Q(x)} té grau \fmm{m} i coeficient principal \fmm{b_m}, en substituir obtenim \fmm{\frac{\infty}{\infty}}. Ens quedem amb els termes de major grau:

\fmdisp{\lim_{x\to\pm\infty}\frac{P(x)}{Q(x)}=\lim_{x\to\pm\infty}\frac{a_nx^n}{b_mx^m}=\begin{cases}\pm\infty&\text{si } n\gt m\ \ (\text{el signe, simplificant})\\[4pt]\dfrac{a_n}{b_m}&\text{si } n=m\\[8pt]0&\text{si } n\lt m\end{cases}}

Les \textbf{arrels} també tenen grau: \fmm{\sqrt{x^3+x}} es comporta com \fmm{\sqrt{x^3}=x^{3/2}} (grau \fmm{\tfrac32}) i \fmm{\sqrt{4x^2+1}} com \fmm{\sqrt{4x^2}=2|x|} (grau 1).
\end{fmdef}

\Needspace{7\baselineskip}
\begin{fmex}{Exemples de quocients}
\begin{enumerate}[label=\textbf{\arabic*.},leftmargin=*]
\item \fmm{\displaystyle\lim_{x\to+\infty}\frac{3x^2-5x+1}{x^3-x}\overset{\frac{\infty}{\infty}}{=}\lim_{x\to+\infty}\frac{3x^2}{x^3}=\lim_{x\to+\infty}\frac{3}{x}=\frac{3}{+\infty}=0} (grau de dalt < grau de baix).
\item \fmm{\displaystyle\lim_{x\to-\infty}\frac{4x^3+x}{2x^3-7}\overset{\frac{\infty}{\infty}}{=}\lim_{x\to-\infty}\frac{4x^3}{2x^3}=\frac42=2} (mateix grau: quocient dels coeficients principals).
\item Grau de dalt més gran: cal mirar el signe a cada costat.

\fmm{\displaystyle\lim_{x\to+\infty}\frac{x^3+1}{2x-x^2}=\lim_{x\to+\infty}\frac{x^3}{-x^2}=\lim_{x\to+\infty}(-x)=-\infty}

\fmm{\displaystyle\lim_{x\to-\infty}\frac{x^3+1}{2x-x^2}=\lim_{x\to-\infty}(-x)=-(-\infty)=+\infty}
\item Amb arrels: \fmm{\displaystyle\lim_{x\to+\infty}\frac{\sqrt{x^3+x-2}}{x+5}=\lim_{x\to+\infty}\frac{x^{3/2}}{x}=\lim_{x\to+\infty}\sqrt{x}=+\infty} i \fmm{\displaystyle\lim_{x\to+\infty}\frac{\sqrt{9x^2+1}}{x-4}=\lim_{x\to+\infty}\frac{3x}{x}=3}.
\end{enumerate}
\end{fmex}

\Needspace{7\baselineskip}
\begin{fmdef}{Arrels quan $x\to-\infty$}
L'arrel quadrada sempre és \textbf{positiva}: \fmm{\sqrt{x^2}=|x|}. Quan \fmm{x\to-\infty}, \fmm{|x|=-x}, de manera que \fmm{\sqrt{a\,x^2+\dots}\approx\sqrt a\,|x|=-\sqrt a\,x}. Una altra manera segura és fer el canvi \fmm{x=-t} amb \fmm{t\to+\infty}.

\begin{enumerate}[label=\textbf{\arabic*.},leftmargin=*]
\item \fmm{\displaystyle\lim_{x\to-\infty}\frac{\sqrt{4x^2+1}}{3x-2}=\lim_{x\to-\infty}\frac{2|x|}{3x}=\lim_{x\to-\infty}\frac{-2x}{3x}=-\frac23} (a \fmm{+\infty} dona \fmm{+\frac23}).
\item \fmm{\displaystyle\lim_{x\to-\infty}\frac{3x+\sqrt{x^2+5}}{2x+1}\overset{x=-t}{=}\lim_{t\to+\infty}\frac{-3t+\sqrt{t^2+5}}{-2t+1}=\lim_{t\to+\infty}\frac{-3t+t}{-2t}=1}.
\item Exponencials de base \fmm{e} en un quocient: dividim per l'exponencial que mana.\\[2ex]
 \fmm{\displaystyle\lim_{x\to+\infty}\frac{3e^{2x}-e^x}{e^{2x}+4}=\lim_{x\to+\infty}\frac{3-e^{-x}}{1+4e^{-2x}}=3}, però \fmm{\displaystyle\lim_{x\to-\infty}\frac{3e^{2x}-e^x}{e^{2x}+4}=\frac{0-0}{0+4}=0}.
\end{enumerate}
\end{fmdef}

\Needspace{18\baselineskip}
\section{Càlcul de límits per comparació d'infinits}\hypertarget{sec3}{}

\Needspace{7\baselineskip}
\begin{fmdef}{Ordre d'un infinit}
Si \fmm{\displaystyle\lim_{x\to+\infty}f(x)=\lim_{x\to+\infty}g(x)=+\infty}, diem que \fmm{f} és un \textbf{infinit d'ordre superior} a \fmm{g} si

\fmdisp{\lim_{x\to+\infty}\frac{f(x)}{g(x)}=+\infty\qquad\left(\text{o, el que és el mateix, }\lim_{x\to+\infty}\frac{g(x)}{f(x)}=0\right)}

És a dir, \fmm{f} creix «molt més de pressa» que \fmm{g}. Si el límit del quocient és un nombre \fmm{k\neq0}, són \textbf{del mateix ordre}.
\end{fmdef}

\Needspace{7\baselineskip}
\begin{fmprop}{Jerarquia d'infinits (quan x $\to$ +$\infty$)}
\begin{center}\fmlevel{logaritmes}{\fmm{\log_a x,\ \ln x}}\quad{\Large\color{fmaccent}$\ll$}\quad\fmlevel{potències}{\fmm{x^p\ \ (p\gt0)}}\quad{\Large\color{fmaccent}$\ll$}\quad\fmlevel{exponencials}{\fmm{a^x\ \ (a\gt1)}}\end{center}

\begin{itemize}
\item Entre dues \textbf{potències}, guanya la d'exponent més gran: \fmm{x^3\gg x^2\gg\sqrt x}.
\item Entre dues \textbf{exponencials}, guanya la de base més gran: \fmm{5^x\gg e^x\gg2^x}.
\item Els coeficients no canvien l'ordre: \fmm{0{,}001\cdot2^x\gg1000\,x^{10}}.
\end{itemize}

\textbf{Conseqüències:} en un \textbf{quocient} \fmm{\frac{\infty}{\infty}}, si el de dalt és d'ordre superior el límit és \fmm{\pm\infty}, i si ho és el de baix, és \fmm{0}. En una \textbf{suma o resta} d'infinits, el resultat el decideix el d'ordre superior.
\end{fmprop}

\Needspace{7\baselineskip}
\begin{fmex}{Exemples}
\begin{enumerate}[label=\textbf{\arabic*.},leftmargin=*]
\item \fmm{\displaystyle\lim_{x\to+\infty}\frac{x^2}{e^x}=0} (l'exponencial de baix és d'ordre superior).
\item \fmm{\displaystyle\lim_{x\to+\infty}\frac{2^x}{x^{10}}=+\infty} (encara que per a \fmm{x} xicotets el denominador és més gran!).
\item \fmm{\displaystyle\lim_{x\to+\infty}\frac{\ln x}{\sqrt x}=0} (el logaritme és d'ordre inferior a qualsevol potència).
\item \fmm{\displaystyle\lim_{x\to+\infty}\left(x^3-2^x\right)=-\infty} (guanya \fmm{-2^x}).
\item Mateix ordre: \fmm{\displaystyle\lim_{x\to+\infty}\frac{e^x+x^5}{3e^x-x}=\lim_{x\to+\infty}\frac{e^x}{3e^x}=\frac13}.
\item \fmm{\displaystyle\lim_{x\to+\infty}\frac{3^x}{2^x}=\lim_{x\to+\infty}\left(\tfrac32\right)^x=+\infty}.
\end{enumerate}
\end{fmex}

\Needspace{18\baselineskip}
\section{Indeterminacions a l'infinit}\hypertarget{sec4}{}

Quan en substituir obtenim una indeterminació, cal \textbf{transformar l'expressió} abans de calcular el límit. Ací veiem les tres que apareixen a la PAU quan \fmm{x\to\pm\infty}.

\Needspace{16\baselineskip}
\subsection*{Indeterminació \fmm{\dfrac{\infty}{\infty}}}

\Needspace{7\baselineskip}
\begin{fmdef}{Dos mètodes equivalents}
\textbf{Mètode 1}: ens quedem amb els \textbf{termes de major grau} del numerador i del denominador i simplifiquem (és el que hem fet a l'apartat 5.2).

\textbf{Mètode 2}: \textbf{dividim} numerador i denominador per la potència de \fmm{x} de major grau de la fracció; tots els termes del tipus \fmm{\frac{k}{x^n}} tendeixen a 0.

\fmdisp{\lim_{x\to+\infty}\frac{3x^2-5x+1}{x^3-x}=\lim_{x\to+\infty}\frac{\frac{3x^2}{x^3}-\frac{5x}{x^3}+\frac{1}{x^3}}{\frac{x^3}{x^3}-\frac{x}{x^3}}=\lim_{x\to+\infty}\frac{\overbrace{\tfrac{3}{x}}^{0}-\overbrace{\tfrac{5}{x^2}}^{0}+\overbrace{\tfrac{1}{x^3}}^{0}}{1-\underbrace{\tfrac{1}{x^2}}_{0}}=\frac01=0}

\end{fmdef}

\Needspace{16\baselineskip}
\subsection*{Indeterminació \fmm{\infty-\infty}}

\Needspace{7\baselineskip}
\begin{fmdef}{Cas 1: resta de fraccions algebraiques}
\textbf{Fem la resta} (amb denominador comú) i calculem el límit de la fracció que en resulta, que sol ser \fmm{\frac{\infty}{\infty}}.
\end{fmdef}

\Needspace{7\baselineskip}
\begin{fmex}{Exemples}
\begin{enumerate}[label=\textbf{\arabic*.},leftmargin=*]
\item \fmm{\displaystyle\lim_{x\to+\infty}\left(\frac{x^2}{x-1}-x\right)}

Substituïm: \fmm{\displaystyle\lim\frac{x^2}{x-1}=\lim\frac{x^2}{x}=+\infty}, i \fmm{\lim x=+\infty\Rightarrow+\infty-\infty\ \ \text{IND.}}

Restem les fraccions: \fmm{\dfrac{x^2}{x-1}-\dfrac{x(x-1)}{x-1}=\dfrac{x^2-x^2+x}{x-1}=\dfrac{x}{x-1}}

Recalculem el límit: \fmm{\displaystyle\lim_{x\to+\infty}\left(\frac{x^2}{x-1}-x\right)=\displaystyle\lim_{x\to+\infty}\frac{x}{x-1}=1} (mateix grau).
\item \fmm{\displaystyle\lim_{x\to+\infty}\left(\frac{3x^2+2}{x+1}-\frac{2x^3-x}{x^2-1}\right)}

\fmm{\displaystyle\lim\frac{3x^2}{x}-\lim\frac{2x^3}{x^2}=\lim3x-\lim2x=+\infty-\infty\ \ \text{IND.}}

Com que \fmm{x^2-1=(x+1)(x-1)}, el denominador comú és \fmm{x^2-1}:
\fmdisp{\frac{(3x^2+2)(x-1)-(2x^3-x)}{x^2-1}=\frac{3x^3-3x^2+2x-2-2x^3+x}{x^2-1}=\frac{x^3-3x^2+3x-2}{x^2-1}}

Per tant: \fmm{\displaystyle\lim_{x\to+\infty}\left(\frac{3x^2+2}{x+1}-\frac{2x^3-x}{x^2-1}\right)=\displaystyle\lim_{x\to+\infty}\frac{x^3-3x^2+3x-2}{x^2-1}=\lim_{x\to+\infty}\frac{x^3}{x^2}=\lim_{x\to+\infty}x=+\infty}
\end{enumerate}
\end{fmex}

\Needspace{7\baselineskip}
\begin{fmdef}{Cas 2: resta amb arrels quadrades}
\textbf{Multipliquem i dividim pel conjugat} i usem \fmm{(A-B)(A+B)=A^2-B^2}, que fa desaparéixer les arrels de dalt:

\fmdisp{\left(\sqrt{A}-B\right)\cdot\frac{\sqrt{A}+B}{\sqrt{A}+B}=\frac{A-B^2}{\sqrt{A}+B}\,\,\,\,\, \text{o bé,}\,\,\,\,\, \left(\sqrt{A}-\sqrt{B}\right)\cdot\frac{\sqrt{A}+\sqrt{B}}{\sqrt{A}+\sqrt{B}}=\frac{A-B}{\sqrt{A}+\sqrt{B}}}
\begin{fmnote}\fmuline{Abans de res}, mira els termes principals: si \fmm{\sqrt{4x^2-2}\approx2x} i li restem \fmm{4x}, ja tenim \fmm{2x-4x\to-\infty} i \textbf{no cal} el conjugat. El conjugat només fa falta quan els termes principals \textbf{es cance\lgem{}en} (com en \fmm{\sqrt{x^2+4x}-x}).\end{fmnote}
\end{fmdef}

\Needspace{7\baselineskip}
\begin{fmex}{Exemples}
\begin{enumerate}[label=\textbf{\arabic*.},leftmargin=*]
\item \fmm{\displaystyle\lim_{x\to+\infty}\left(\sqrt{x^2+4x}-x\right)=\sqrt{+\infty}-\infty=+\infty-\infty\ \ \text{IND.}} (els termes \fmm{\sqrt{x^2}=x} i \fmm{x} es cance\lgem{}en).\\[2ex]

Multipliquem i dividim pel conjugat:
\fmdisp{\left(\sqrt{x^2+4x}-x\right)\frac{\sqrt{x^2+4x}+x}{\sqrt{x^2+4x}+x}=\frac{x^2+4x-x^2}{\sqrt{x^2+4x}+x}=\frac{4x}{\sqrt{x^2+4x}+x}}

I tenim: \fmm{\displaystyle\lim_{x\to+\infty}\left(\sqrt{x^2+4x}-x\right)=\displaystyle\lim_{x\to+\infty}\frac{4x}{\sqrt{x^2+4x}+x}=\lim_{x\to+\infty}\frac{4x}{x+x}=\frac{4}{2}=2}
\item \fmm{\displaystyle\lim_{x\to+\infty}\left(\sqrt{x^2+x}-\sqrt{x^2-x}\right)=+\infty-\infty\ \ \text{IND.}} Multipliquem i dividim per \fmm{\sqrt{x^2+x}+\sqrt{x^2-x}}:
\fmdisp{\frac{(x^2+x)-(x^2-x)}{\sqrt{x^2+x}+\sqrt{x^2-x}}=\frac{2x}{\sqrt{x^2+x}+\sqrt{x^2-x}}}

Aleshores: \fmm{\displaystyle\lim_{x\to+\infty}\left(\sqrt{x^2+x}-\sqrt{x^2-x}\right)=\displaystyle\lim_{x\to+\infty}\frac{2x}{x+x}=1}
\end{enumerate}
\end{fmex}

\Needspace{16\baselineskip}
\subsection*{Indeterminació \fmm{1^{\infty}}}

\Needspace{7\baselineskip}
\begin{fmprop}{Fórmula per a la indeterminació 1\textsuperscript{$\infty$}}
Si \fmm{\displaystyle\lim f(x)=1} i \fmm{\displaystyle\lim g(x)=\pm\infty}, tenim \fmm{1^{\infty}}. Llavors:

\fmdisp{\lim_{x\to\pm\infty}f(x)^{g(x)}=e^{\lim\limits_{x\to\pm\infty}\big(f(x)-1\big)\cdot g(x)}}

\begin{enumerate}
\item Restem 1 a la base: \fmm{f(x)-1} (normalment, fent la resta de fraccions).
\item Multipliquem per l'exponent: \fmm{(f(x)-1)\cdot g(x)}.
\item Calculem el límit \fmm{\lambda} d'aquest producte; el resultat és \fmm{e^{\lambda}} (amb \fmm{e^{+\infty}=+\infty} i \fmm{e^{-\infty}=0}).
\end{enumerate}

\begin{fmnote}\fmuline{Compte}: si la base tendeix a un nombre \fmm{k\neq1}, \textbf{no} hi ha indeterminació: \fmm{\displaystyle\lim_{x\to+\infty}\left(\frac{2x+1}{x+3}\right)^{x}=2^{+\infty}=+\infty}, \fmm{\ \displaystyle\lim_{x\to+\infty}\left(\frac{x+1}{3x}\right)^{x^2}=\left(\tfrac13\right)^{+\infty}=0}.\end{fmnote}
\end{fmprop}

\Needspace{7\baselineskip}
\begin{fmex}{Exemples}
\begin{enumerate}[label=\textbf{\arabic*.},leftmargin=*]
\item \fmm{\displaystyle\lim_{x\to-\infty}\left(\frac{x-3}{x-1}\right)^{2x-1}}

\fmm{\displaystyle\lim_{x\to-\infty}\frac{x-3}{x-1}=1} i \fmm{\displaystyle\lim_{x\to-\infty}(2x-1)=-\infty\Rightarrow1^{-\infty}\ \ \text{IND.}}

1.- \fmm{\ \dfrac{x-3}{x-1}-1=\dfrac{x-3-(x-1)}{x-1}=\dfrac{-2}{x-1}}

2.- \fmm{\ \dfrac{-2}{x-1}\cdot(2x-1)=\dfrac{-4x+2}{x-1}}

3.- \fmm{\ \displaystyle\lim_{x\to-\infty}\frac{-4x+2}{x-1}=-4\Rightarrow\lim_{x\to-\infty}\left(\frac{x-3}{x-1}\right)^{2x-1}=e^{-4}=\frac{1}{e^4}}
\item \fmm{\displaystyle\lim_{x\to+\infty}\left(\frac{2x+1}{2x-3}\right)^{x}}

\fmm{1^{+\infty}\ \ \text{IND.}} \fmm{\quad\dfrac{2x+1}{2x-3}-1=\dfrac{4}{2x-3}}

\fmm{\displaystyle\lim_{x\to+\infty}\frac{4}{2x-3}\cdot x=\lim_{x\to+\infty}\frac{4x}{2x-3}=2\Rightarrow} el límit és \fmm{e^2}
\item \fmm{\displaystyle\lim_{x\to+\infty}\left(1+\frac{3}{x}\right)^{x}=e^{\lim\frac3x\cdot x}=e^{3}}. I, en canvi, \fmm{\displaystyle\lim_{x\to+\infty}\left(\frac{x^2+1}{x^2}\right)^{x}=e^{\lim\frac{1}{x^2}\cdot x}=e^{0}=1}.
\end{enumerate}
\end{fmex}

\Needspace{7\baselineskip}
\begin{fmdef}{D'on ix la fórmula?}
El nombre \fmm{e} es defineix com \fmm{\displaystyle e=\lim_{x\to\pm\infty}\left(1+\frac1x\right)^{x}\approx2{,}71828}. Si \fmm{f\to1} i \fmm{g\to\pm\infty}, escrivim \fmm{f=1+\frac{1}{h}} amb \fmm{h=\frac{1}{f-1}\to\pm\infty}:

\fmdisp{f^{\,g}=\left[\left(1+\frac1h\right)^{h}\right]^{\frac{g}{h}}\longrightarrow e^{\lim\frac gh}=e^{\lim\,(f-1)\,g}}

La fórmula també val quan \fmm{x} tendeix a un nombre \fmm{a} (apartat 5.6).
\end{fmdef}

\Needspace{16\baselineskip}
\subsection*{Indeterminació \fmm{0\cdot\infty}}

\Needspace{7\baselineskip}
\begin{fmdef}{Mètode}
Transformem el producte en un quocient per arribar a \fmm{\frac{\infty}{\infty}} o \fmm{\frac00}: \fmm{\ f\cdot g=\dfrac{f}{1/g}}. Sovint n'hi ha prou d'\textbf{operar} (ajuntar fraccions, multiplicar) o de comparar infinits.
\end{fmdef}

\Needspace{7\baselineskip}
\begin{fmex}{Exemples}
\begin{enumerate}[label=\textbf{\arabic*.},leftmargin=*]
\item \fmm{\displaystyle\lim_{x\to+\infty}\frac{3}{x^2+1}\cdot(2x^2-x)\overset{0\cdot\infty}{=}\lim_{x\to+\infty}\frac{6x^2-3x}{x^2+1}=6}.
\item \fmm{\displaystyle\lim_{x\to+\infty}x\left(\sqrt{x^2+4}-x\right)\overset{\infty\cdot0}{=}\lim_{x\to+\infty}\frac{x\,(x^2+4-x^2)}{\sqrt{x^2+4}+x}=\lim_{x\to+\infty}\frac{4x}{2x}=2}.
\item \fmm{\displaystyle\lim_{x\to+\infty}x^3e^{-x}\overset{\infty\cdot0}{=}\lim_{x\to+\infty}\frac{x^3}{e^x}=0} (l'exponencial és d'ordre superior).
\item \fmm{\displaystyle\lim_{x\to+\infty}x\ln\!\left(1+\frac2x\right)=\lim_{x\to+\infty}\ln\left(1+\frac2x\right)^{x}=\ln e^{2}=2} (hem passat l'exponent dins del logaritme i hem aplicat \fmm{1^\infty}).
\end{enumerate}
\end{fmex}

\Needspace{16\baselineskip}
\subsection*{Indeterminacions \fmm{\infty^0} i \fmm{0^0}}

\Needspace{7\baselineskip}
\begin{fmdef}{Mètode: prenem logaritmes}
Si \fmm{L=\lim f^{\,g}}, llavors \fmm{\ln L=\lim g\cdot\ln f}, que sol ser una indeterminació \fmm{0\cdot\infty}. En acabar, desfem el logaritme: \fmm{L=e^{\ln L}}.

\begin{enumerate}[label=\textbf{\arabic*.},leftmargin=*]
\item \fmm{\displaystyle L=\lim_{x\to+\infty}x^{1/x}\ (\infty^0)}: \fmm{\ \displaystyle\ln L=\lim_{x\to+\infty}\frac{\ln x}{x}=0} (\fmm{x} és d'ordre superior a \fmm{\ln x}), així que \fmm{L=e^0=1}.
\item \fmm{\displaystyle L=\lim_{x\to+\infty}\left(x^2+1\right)^{\frac{1}{\ln x}}\ (\infty^0)}: \fmm{\ \displaystyle\ln L=\lim_{x\to+\infty}\frac{\ln(x^2+1)}{\ln x}=\lim_{x\to+\infty}\frac{\ln x^2}{\ln x}=\lim_{x\to+\infty}\frac{2\ln x}{\ln x}=2\Rightarrow L=e^2}.
\end{enumerate}

\begin{fmnote}\textbf{Resum}: \fmm{\frac{\infty}{\infty}} $\to$ termes dominants; \fmm{\infty-\infty} $\to$ operar o conjugat; \fmm{0\cdot\infty} $\to$ quocient;\\[2ex]
 \fmm{1^\infty} $\to$ \fmm{e^{\lim(f-1)g}}; \fmm{\infty^0}, \fmm{0^0} $\to$ logaritmes.\end{fmnote}
\end{fmdef}

\Needspace{18\baselineskip}
\section{Límits en un nombre: càlcul gràfic i analític}\hypertarget{sec5}{}

\Needspace{7\baselineskip}
\begin{fmdef}{Límit en un punt i límits laterals}
\fmm{\displaystyle\lim_{x\to a}f(x)=L} vol dir que, quan \fmm{x} s'acosta a \fmm{a} (sense arribar a ser \fmm{a}), \fmm{f(x)} s'acosta a \fmm{L}. Ens podem acostar a \fmm{a} per dos costats:

\begin{itemize}
\item \textbf{Límit per l'esquerra} \fmm{\displaystyle\lim_{x\to a^-}f(x)}: amb valors \fmm{x\lt a} (per exemple \fmm{a-0{,}1,\ a-0{,}01,\dots}).
\item \textbf{Límit per la dreta} \fmm{\displaystyle\lim_{x\to a^+}f(x)}: amb valors \fmm{x\gt a}.
\end{itemize}

\fmdisp{\exists\lim_{x\to a}f(x)=L\iff\lim_{x\to a^-}f(x)=\lim_{x\to a^+}f(x)=L}

Si algun límit lateral és \fmm{\pm\infty}, la recta \fmm{x=a} és una \textbf{asímptota vertical}. Observa que el límit en \fmm{a} \textbf{no depén} del valor \fmm{f(a)}: pot ser que \fmm{f(a)} ni existisca.
\end{fmdef}

\Needspace{7\baselineskip}
\begin{fmex}{Càlcul gràfic}
\begin{center}\begin{tikzpicture}\begin{axis}[fmaxis,width=8cm,height=6cm,xmin=-4,xmax=6,ymin=-4,ymax=6]\draw[fmasy] ({axis cs:-1,0}|-{rel axis cs:0,0}) -- ({axis cs:-1,0}|-{rel axis cs:0,1});\addplot[fmcurve,domain=-4:-1.1,samples=300]{-1/(x+1)};\addplot[fmcurve,domain=-1:2,samples=160]{x};\addplot[fmcurve,domain=2:6,samples=160]{4-(x-2)^2/2};\node[fmopen] at (axis cs:-1,-1) {};\node[fmopen] at (axis cs:1,1) {};\node[fmdot] at (axis cs:1,3) {};\node[fmopen] at (axis cs:2,2) {};\node[fmdot] at (axis cs:2,4) {};\end{axis}\end{tikzpicture}\end{center}

Llegim a la gràfica de \fmm{f} (un punt buit $\circ$ vol dir que el punt \textbf{no} és de la gràfica):

\begin{center}\renewcommand{\arraystretch}{1.5}\small\begin{tabularx}{\linewidth}{>{\raggedright\arraybackslash}p{3.2cm}>{\raggedright\arraybackslash}X>{\raggedright\arraybackslash}X>{\raggedright\arraybackslash}X>{\raggedright\arraybackslash}X}\toprule
\textbf{$x\to$} & \textbf{Per l'esquerra} & \textbf{Per la dreta} & \textbf{Límit} & \textbf{$f(a)$}\\\midrule
$-1$ & $+\infty$ & $-1$ & $\nexists$ & $\nexists$\\[3pt]
$1$ & $1$ & $1$ & $1$ & $3$\\[3pt]
$2$ & $2$ & $4$ & $\nexists$ & $4$\\[3pt]
\bottomrule\end{tabularx}\end{center}

A més, \fmm{\displaystyle\lim_{x\to-\infty}f(x)=0} (asímptota horitzontal \fmm{y=0}) i \fmm{\displaystyle\lim_{x\to+\infty}f(x)=-\infty}.
\end{fmex}

\Needspace{7\baselineskip}
\begin{fmdef}{Càlcul analític}
\begin{enumerate}
\item \textbf{Substituïm} \fmm{x=a}. Si obtenim un nombre, aquest és el límit: \fmm{\displaystyle\lim_{x\to2}\left(x^2-3x+5\right)=4-6+5=3}.
\item Si obtenim \fmm{\dfrac{k}{0}} amb \fmm{k\neq0}, el límit és infinit: calculem els \textbf{límits laterals} per saber el signe (substituint valors propers, com \fmm{a-0{,}1} i \fmm{a+0{,}1}).
\item Si obtenim \fmm{\dfrac00}, és una indeterminació: la resoldrem a l'apartat 5.6.
\item En una \textbf{funció a trossos}, en els punts on canvia l'expressió calculem cada límit lateral amb el tros que li correspon.
\end{enumerate}
\end{fmdef}

\Needspace{7\baselineskip}
\begin{fmex}{Exemples}
\begin{enumerate}[label=\textbf{\arabic*.},leftmargin=*]
\item Funció a trossos: \fmm{f(x)=\begin{cases}2x^2-6&\text{si } x\leq1\\[2pt]\dfrac{2x}{x-3}&\text{si } x\gt1\end{cases}}. Calculem \fmm{\displaystyle\lim_{x\to1}f(x)}.

\fmm{\displaystyle\lim_{x\to1^-}f(x)=\lim_{x\to1}\left(2x^2-6\right)=2-6=-4}

\fmm{\displaystyle\lim_{x\to1^+}f(x)=\lim_{x\to1}\frac{2x}{x-3}=\frac{2}{-2}=-1}

Els límits laterals són diferents: \fmm{\nexists\displaystyle\lim_{x\to1}f(x)}.
\item \fmm{\displaystyle\lim_{x\to1}\frac{4x^2}{x^2-1}}

Substituïm: \fmm{\dfrac{4}{0}\Rightarrow} límit infinit; estudiem el signe a cada costat.

\fmm{x=0{,}9:\ \dfrac{4\cdot0{,}81}{0{,}81-1}=\dfrac{+}{-}\lt0\Rightarrow\displaystyle\lim_{x\to1^-}f(x)=-\infty}

\fmm{x=1{,}1:\ \dfrac{4\cdot1{,}21}{1{,}21-1}=\dfrac{+}{+}\gt0\Rightarrow\displaystyle\lim_{x\to1^+}f(x)=+\infty}

\fmm{\nexists\displaystyle\lim_{x\to1}f(x)} i \fmm{x=1} és una asímptota vertical.
\item \fmm{\displaystyle\lim_{x\to2}\frac{3}{(x-2)^2}=\frac30}: el denominador és sempre positiu, així que els dos laterals són \fmm{+\infty} i \fmm{\displaystyle\lim_{x\to2}\frac{3}{(x-2)^2}=+\infty}.
\item \fmm{\displaystyle\lim_{x\to3^+}\log(x-3)=\log(0^+)=-\infty}; per l'esquerra no existeix (el logaritme d'un nombre negatiu no està definit).
\end{enumerate}
\end{fmex}

\Needspace{18\baselineskip}
\section{Indeterminacions en un nombre: $\frac00$, $\infty-\infty$ i $1^\infty$}\hypertarget{sec6}{}

\Needspace{7\baselineskip}
\begin{fmdef}{Com es resol la indeterminació 0/0}
Si \fmm{P(a)=0} i \fmm{Q(a)=0}, pel teorema del factor \textbf{tots dos polinomis són divisibles per \fmm{(x-a)}}. Per tant:

\begin{enumerate}
\item \textbf{Factoritzem} numerador i denominador: traient factor comú, amb les identitats notables, resolent l'equació de 2n grau o amb Ruffini.
\item \textbf{Simplifiquem} el factor \fmm{(x-a)} (podem fer-ho perquè \fmm{x\neq a} quan calculem el límit).
\item \textbf{Tornem a substituir}. Si torna a eixir \fmm{\frac00}, repetim; si ix \fmm{\frac{k}{0}}, estudiem els límits laterals.
\end{enumerate}
\end{fmdef}

\Needspace{7\baselineskip}
\begin{fmex}{Exemples}
\begin{enumerate}[label=\textbf{\arabic*.},leftmargin=*]
\item \fmm{\displaystyle\lim_{x\to2}\frac{x^2-x-2}{x^2-4}}

Substituïm: \fmm{\dfrac{4-2-2}{4-4}=\dfrac00\ \ \text{IND.}}

Factoritzem (\fmm{x^2-x-2=0\Rightarrow x=2,\ x=-1}; \fmm{\ x^2-4} és una diferència de quadrats):
\fmdisp{\lim_{x\to2}\frac{\cancel{(x-2)}(x+1)}{\cancel{(x-2)}(x+2)}=\lim_{x\to2}\frac{x+1}{x+2}=\frac34}

\item \fmm{\displaystyle\lim_{x\to1}\frac{x^3-3x+2}{x^2-1}}

\fmm{\dfrac{1-3+2}{1-1}=\dfrac00\ \ \text{IND.}} Dividim el numerador per \fmm{(x-1)} amb Ruffini:

\begin{center}\begin{tabular}{r|rrrr}
 & $1$ & $0$ & $-3$ & $2$\\
$1$ &  & $1$ & $1$ & $-2$\\\hline
 & $1$ & $1$ & $-2$ & \fbox{$0$}\\
\end{tabular}\end{center}

\fmm{x^3-3x+2=(x-1)(x^2+x-2)=(x-1)(x-1)(x+2)}

\fmdisp{\lim_{x\to1}\frac{(x-1)^2(x+2)}{(x-1)(x+1)}=\lim_{x\to1}\frac{(x-1)(x+2)}{x+1}=\frac{0\cdot3}{2}=0}

\item \fmm{\displaystyle\lim_{x\to-1}\frac{x^2+x}{x^2+2x+1}}

\fmm{\dfrac00\ \ \text{IND.}} \fmm{\quad\displaystyle\lim_{x\to-1}\frac{x\cancel{(x+1)}}{(x+1)^{\cancel{2}}}=\lim_{x\to-1}\frac{x}{x+1}=\frac{-1}{0}}

Ara ix \fmm{\frac{k}{0}}: laterals. \fmm{\ x=-1{,}1:\ \frac{-}{-}\gt0\Rightarrow\displaystyle\lim_{x\to-1^-}=+\infty}; \fmm{\ x=-0{,}9:\ \frac{-}{+}\lt0\Rightarrow\displaystyle\lim_{x\to-1^+}=-\infty}. No existeix el límit.
\item Amb arrels, pel conjugat: \fmm{\displaystyle\lim_{x\to4}\frac{\sqrt x-2}{x-4}=\lim_{x\to4}\frac{(\sqrt x-2)(\sqrt x+2)}{(x-4)(\sqrt x+2)}=\lim_{x\to4}\frac{\cancel{x-4}}{\cancel{(x-4)}(\sqrt x+2)}=\frac14}
\end{enumerate}
\end{fmex}

\Needspace{16\baselineskip}
\subsection*{Arrels: multipliquem pel conjugat}

\Needspace{7\baselineskip}
\begin{fmex}{Exemples}
\begin{enumerate}[label=\textbf{\arabic*.},leftmargin=*]
\item \fmm{\displaystyle\lim_{x\to3}\frac{\sqrt{x+1}-\sqrt{2x-2}}{x^2-9}}

\fmm{\dfrac{2-2}{0}=\dfrac00\ \ \text{IND.}} Multipliquem i dividim pel conjugat del numerador:
\fmdisp{\frac{(x+1)-(2x-2)}{(x^2-9)\left(\sqrt{x+1}+\sqrt{2x-2}\right)}=\frac{-(x-3)}{(x-3)(x+3)\left(\sqrt{x+1}+\sqrt{2x-2}\right)}}

\fmm{\displaystyle\lim_{x\to3}\frac{-1}{(x+3)\left(\sqrt{x+1}+\sqrt{2x-2}\right)}=\frac{-1}{6\cdot4}=-\frac{1}{24}}
\item Arrels al numerador i al denominador: dos conjugats.
\fmdisp{\lim_{x\to0}\frac{\sqrt{1+x}-1}{\sqrt{4+x}-2}=\lim_{x\to0}\frac{x\left(\sqrt{4+x}+2\right)}{x\left(\sqrt{1+x}+1\right)}=\frac{4}{2}=2}

\end{enumerate}
\end{fmex}

\Needspace{2\baselineskip}
\subsection*{Indeterminacions \fmm{\infty-\infty} i \fmm{1^\infty} en un nombre}

\Needspace{7\baselineskip}
\begin{fmdef}{Mètode}
\begin{itemize}
\item \fmm{\infty-\infty}: fem la resta de fraccions; normalment ix \fmm{\frac00} o \fmm{\frac k0}.
\item \fmm{1^\infty}: la mateixa fórmula que a l'infinit, \fmm{\displaystyle\lim_{x\to a}f(x)^{g(x)}=e^{\lim\limits_{x\to a}(f(x)-1)\,g(x)}}. Compte: si l'exponent tendeix a \fmm{+\infty} per un costat i a \fmm{-\infty} per l'altre, cal calcular els límits laterals.
\end{itemize}
\end{fmdef}

\Needspace{7\baselineskip}
\begin{fmex}{Exemples}
\begin{enumerate}[label=\textbf{\arabic*.},leftmargin=*]
\item \fmm{\displaystyle\lim_{x\to1}\left(\frac{1}{x-1}-\frac{2}{x^2-1}\right)\overset{\infty-\infty}{=}\lim_{x\to1}\frac{x+1-2}{(x-1)(x+1)}=\lim_{x\to1}\frac{\cancel{x-1}}{\cancel{(x-1)}(x+1)}=\frac12}
\item \fmm{\displaystyle\lim_{x\to0}\left(1+3x\right)^{\frac2x}\overset{1^\infty}{=}e^{\lim\limits_{x\to0}3x\cdot\frac2x}=e^{6}}
\item \fmm{\displaystyle\lim_{x\to2}\left(\frac{x+2}{2x}\right)^{\frac{1}{x-2}}}

\fmm{1^{\frac10}}: la base tendeix a 1 i l'exponent a \fmm{\pm\infty}. \fmm{\quad\dfrac{x+2}{2x}-1=\dfrac{2-x}{2x}}

\fmm{\displaystyle\lim_{x\to2}\frac{-(x-2)}{2x}\cdot\frac{1}{x-2}=\lim_{x\to2}\frac{-1}{2x}=-\frac14\Rightarrow} el límit és \fmm{e^{-1/4}=\dfrac{1}{\sqrt[4]{e}}} (pels dos costats).
\item \fmm{\displaystyle\lim_{x\to0}\left(1+x^2\right)^{\frac{1}{x}}=e^{\lim\limits_{x\to0}x^2\cdot\frac1x}=e^{0}=1}.
\end{enumerate}
\end{fmex}

\Needspace{18\baselineskip}
\section{Límits trigonomètrics, exponencials i logarítmics. Infinitèsims equivalents}\hypertarget{sec7}{}

\Needspace{7\baselineskip}
\begin{fmprop}{Límit fonamental de la trigonometria}

\fmdisp{\lim_{x\to0}\frac{\sin x}{x}=1\qquad(x\ \text{en radians})}

\par\noindent \begin{minipage}[t]{0.48\linewidth}\vspace{0pt}\raggedright \textbf{Per què?} En la circumferència de radi 1, per a \fmm{0\lt x\lt\frac\pi2}, comparem les àrees del triangle \fmm{OAP}, del sector circular \fmm{OAP} i del triangle \fmm{OAT}:

\fmdisp{\frac{\sin x}{2}\lt\frac{x}{2}\lt\frac{\tan x}{2}\ \Rightarrow\ 1\lt\frac{x}{\sin x}\lt\frac{1}{\cos x}\ \Rightarrow\ \cos x\lt\frac{\sin x}{x}\lt1}

Quan \fmm{x\to0}, \fmm{\cos x\to1}: la funció \fmm{\frac{\sin x}{x}} queda «atrapada» entre dues funcions que tendeixen a 1 (\textbf{regla del sandvitx}). Com que \fmm{\frac{\sin x}{x}} és parella, també val per l'esquerra.\end{minipage}\hfill\begin{minipage}[t]{0.48\linewidth}\vspace{0pt}\raggedright \begin{center}\begin{tikzpicture}[x=0.034cm,y=0.034cm,line cap=round,font=\small]
  \fill[fmaccent,fill opacity=0.18] (0,0)--(120,0) arc[start angle=0,end angle=34.4,radius=120]--cycle;
  \draw[fmmuted!70,line width=0.9pt] (120,0) arc[start angle=0,end angle=90,radius=120];
  \draw[line width=0.8pt] (0,0)--(132,0); \draw[line width=0.9pt] (0,0)--(120,82.1);
  \draw[fmaccent,line width=1.8pt] (99,67.7)--(99,0);
  \draw[fmex,line width=1.8pt] (120,0)--(120,82.1);
  \foreach \p in {(0,0),(99,67.7),(120,0),(120,82.1)} \fill \p circle[radius=2.6];
  \node[below left] at (0,0) {$O$}; \node[below] at (120,0) {$A$}; \node[above left] at (99,67.7) {$P$}; \node[right] at (120,82.1) {$T$};
  \node[fmaccent,left] at (99,30) {$\sin x$}; \node[fmex,right] at (120,40) {$\tan x$}; \node at (110,29) {$x$};
\end{tikzpicture}\end{center}\end{minipage}\par\medskip
\end{fmprop}

\Needspace{7\baselineskip}
\begin{fmdef}{Conseqüències}
\begin{itemize}
\item \fmm{\displaystyle\lim_{x\to0}\frac{\sin(kx)}{x}=\lim_{x\to0}k\,\frac{\sin(kx)}{kx}=k}, \fmm{\quad\displaystyle\lim_{x\to0}\frac{\tan x}{x}=\lim_{x\to0}\frac{\sin x}{x}\cdot\frac{1}{\cos x}=1}.
\item \fmm{\displaystyle\lim_{x\to0}\frac{1-\cos x}{x^2}=\lim_{x\to0}\frac{(1-\cos x)(1+\cos x)}{x^2(1+\cos x)}=\lim_{x\to0}\left(\frac{\sin x}{x}\right)^2\frac{1}{1+\cos x}=\frac12}.
\item Del nombre \fmm{e} (apartat 5.4) ixen dos límits més: \fmm{\ln(1+x)^{1/x}\to\ln e=1} i, amb el canvi \fmm{t=e^x-1}, el de l'exponencial:
\fmdisp{\lim_{x\to0}\frac{\ln(1+x)}{x}=1\qquad\qquad\lim_{x\to0}\frac{e^x-1}{x}=1}

\end{itemize}
\end{fmdef}

\Needspace{16\baselineskip}
\subsection*{Infinitèsims equivalents}

\Needspace{7\baselineskip}
\begin{fmdef}{Definició i taula}
Una funció és un \textbf{infinitèsim} en \fmm{x=a} si \fmm{\lim_{x\to a}f(x)=0}. Dos infinitèsims són \textbf{equivalents}, \fmm{f\sim g}, si \fmm{\lim_{x\to a}\frac{f(x)}{g(x)}=1}. Quan \fmm{u(x)\to0}:

\begin{center}\renewcommand{\arraystretch}{1.5}\small\begin{tabularx}{\linewidth}{>{\raggedright\arraybackslash}p{5.2cm}>{\raggedright\arraybackslash}X>{\raggedright\arraybackslash}X}\toprule
$\sin u\sim u$ & $\tan u\sim u$ & $1-\cos u\sim\dfrac{u^2}{2}$\\[3pt]
$e^{u}-1\sim u$ & $\ln(1+u)\sim u$ & $a^{u}-1\sim u\ln a$\\[3pt]
\bottomrule\end{tabularx}\end{center}

\textbf{Regla}: en un \textbf{producte o quocient} podem substituir un infinitèsim per un altre d'equivalent, i el límit no canvia.

\begin{fmnote}\fmuline{Compte}: \textbf{no} es pot fer en sumes ni restes. Per exemple, \fmm{\frac{\tan x-\sin x}{x^3}} no és \fmm{\frac{x-x}{x^3}=0}: cal escriure \fmm{\tan x-\sin x=\tan x\,(1-\cos x)\sim x\cdot\frac{x^2}{2}}, i el límit és \fmm{\frac12}.\end{fmnote}
\end{fmdef}

\Needspace{7\baselineskip}
\begin{fmex}{Exemples}
\begin{enumerate}[label=\textbf{\arabic*.},leftmargin=*]
\item \fmm{\displaystyle\lim_{x\to0}\frac{\sin 3x}{\tan 5x}\overset{\frac00}{=}\lim_{x\to0}\frac{3x}{5x}=\frac35}
\item \fmm{\displaystyle\lim_{x\to0}\frac{1-\cos 2x}{x\sin x}\overset{\frac00}{=}\lim_{x\to0}\frac{\frac{(2x)^2}{2}}{x\cdot x}=\lim_{x\to0}\frac{2x^2}{x^2}=2}
\item \fmm{\displaystyle\lim_{x\to0}\frac{e^{2x}-1}{\ln(1+3x)}\overset{\frac00}{=}\lim_{x\to0}\frac{2x}{3x}=\frac23}
\item Quan el punt no és el 0, fem un canvi de variable perquè ho siga:

\fmm{\displaystyle\lim_{x\to\pi}\frac{\sin x}{x-\pi}\overset{\frac00}{=}} amb \fmm{t=x-\pi\to0}: \fmm{\ \sin x=\sin(t+\pi)=-\sin t}

\fmm{\displaystyle\lim_{t\to0}\frac{-\sin t}{t}=-1}
\item \fmm{\displaystyle\lim_{x\to1}\frac{\ln x}{x^2-1}}: \fmm{\ t=x-1\to0}, \fmm{\ \ln x=\ln(1+t)\sim t}: \fmm{\ \displaystyle\lim_{t\to0}\frac{t}{t(t+2)}=\frac12}
\item \textbf{Fitada per infinitèsim}: si \fmm{|g(x)|\leq M} i \fmm{f(x)\to0}, aleshores \fmm{f(x)\,g(x)\to0}. Així, \fmm{\displaystyle\lim_{x\to+\infty}\frac{\sin x}{x}=0} i \fmm{\displaystyle\lim_{x\to0}x^2\cos\frac1x=0}. En canvi, \fmm{\displaystyle\lim_{x\to+\infty}x\sin\frac1x\overset{t=1/x}{=}\lim_{t\to0^+}\frac{\sin t}{t}=1}.
\end{enumerate}
\end{fmex}

\Needspace{18\baselineskip}
\section{Límits que depenen d'un paràmetre}\hypertarget{sec8}{}

\Needspace{7\baselineskip}
\begin{fmdef}{Mètode}
\begin{enumerate}
\item Calculem el límit \textbf{com si el paràmetre fóra un nombre}, amb les tècniques dels apartats anteriors. El resultat serà una expressió amb el paràmetre.
\item \textbf{Imposem la condició} de l'enunciat (que valga un nombre donat, que siga finit\dots{}) i resolem l'equació.
\item Si el comportament canvia segons el valor del paràmetre, \textbf{discutim} els casos per separat.
\end{enumerate}

\begin{fmnote}\fmuline{Clau} per a \fmm{\displaystyle\lim_{x\to a}\frac{P(x)}{Q(x)}} amb \fmm{Q(a)=0}: perquè el límit siga \textbf{finit}, cal que també \fmm{P(a)=0} (si no, seria \fmm{\frac{k}{0}=\infty}).\end{fmnote}
\end{fmdef}

\Needspace{7\baselineskip}
\begin{fmex}{Exemples}
\begin{enumerate}[label=\textbf{\arabic*.},leftmargin=*]
\item Troba \fmm{a} perquè \fmm{\displaystyle\lim_{x\to+\infty}\left(\sqrt{x^2+ax}-x\right)=3}.

\fmm{\displaystyle\lim_{x\to+\infty}\frac{x^2+ax-x^2}{\sqrt{x^2+ax}+x}=\lim_{x\to+\infty}\frac{ax}{2x}=\frac a2}

\fmm{\dfrac a2=3\Rightarrow a=6}
\item Troba \fmm{a} perquè \fmm{\displaystyle\lim_{x\to+\infty}\left(\frac{x+a}{x-1}\right)^{x}=e^{3}}.

\fmm{1^\infty}: \fmm{\ \dfrac{x+a}{x-1}-1=\dfrac{a+1}{x-1}\Rightarrow\displaystyle\lim_{x\to+\infty}\frac{(a+1)x}{x-1}=a+1\Rightarrow} el límit és \fmm{e^{a+1}}

\fmm{e^{a+1}=e^{3}\Rightarrow a+1=3\Rightarrow a=2}
\item Troba \fmm{a} perquè el límit siga finit i calcula'l: \fmm{\displaystyle\lim_{x\to2}\frac{x^2+ax-6}{x^2-4}}.

El denominador s'anu\lgem{}a en \fmm{x=2}. El numerador també ho ha de fer: \fmm{4+2a-6=0\Rightarrow a=1}

\fmm{\displaystyle\lim_{x\to2}\frac{x^2+x-6}{x^2-4}=\lim_{x\to2}\frac{\cancel{(x-2)}(x+3)}{\cancel{(x-2)}(x+2)}=\frac54}. Per a \fmm{a\neq1} ix \fmm{\frac{k}{0}}: el límit no és finit.
\item Dos paràmetres: troba \fmm{a} i \fmm{b} perquè \fmm{\displaystyle\lim_{x\to1}\frac{x^2+ax+b}{x-1}=3}.

Ha de ser \fmm{\frac00}: \fmm{1+a+b=0\Rightarrow b=-1-a}. Aleshores \fmm{x^2+ax-1-a=(x-1)(x+1+a)} (Ruffini).

\fmm{\displaystyle\lim_{x\to1}(x+1+a)=2+a=3\Rightarrow a=1,\ b=-2}
\item Amb infinitèsims: troba \fmm{a\gt0} perquè \fmm{\displaystyle\lim_{x\to0}\frac{1-\cos(ax)}{x\sin x}=8}.

\fmm{\displaystyle\lim_{x\to0}\frac{\frac{a^2x^2}{2}}{x\cdot x}=\frac{a^2}{2}=8\Rightarrow a^2=16\Rightarrow a=4}
\item \textbf{Discussió}: calcula, segons els valors de \fmm{a}, \fmm{\displaystyle\lim_{x\to+\infty}\left(\sqrt{x^2+1}-ax\right)}.

\fmm{\sqrt{x^2+1}-ax\approx x-ax=(1-a)x}. Si \fmm{a\lt1}, el límit és \fmm{+\infty}; si \fmm{a\gt1}, és \fmm{-\infty}.

Si \fmm{a=1}, és \fmm{\infty-\infty}: \fmm{\displaystyle\lim_{x\to+\infty}\frac{1}{\sqrt{x^2+1}+x}=0}.
\end{enumerate}
\end{fmex}

\Needspace{18\baselineskip}
\section{Continuïtat d'una funció en un punt}\hypertarget{sec9}{}

\Needspace{7\baselineskip}
\begin{fmdef}{Definició: funció contínua en un punt}
Intuïtivament, \fmm{f} és contínua en \fmm{x=a} si la seua gràfica es pot dibuixar al voltant de \fmm{a} \textbf{sense alçar el llapis}. Formalment, \fmm{f} és \textbf{contínua en \fmm{x=a}} si es compleixen les tres condicions:

\begin{enumerate}
\item \fmm{\exists f(a)} (la funció està definida en \fmm{a}).
\item \fmm{\exists\displaystyle\lim_{x\to a}f(x)} i és un nombre (els dos límits laterals coincideixen).
\item \fmm{\displaystyle\lim_{x\to a}f(x)=f(a)}.
\end{enumerate}

Si falla alguna condició, \fmm{f} és \textbf{discontínua} en \fmm{a}. Una funció és contínua en un interval si ho és en tots els punts de l'interval.
\end{fmdef}

\Needspace{7\baselineskip}
\begin{fmprop}{Tipus de discontinuïtat}
\par\noindent \begin{minipage}[t]{0.235\linewidth}\centering {\centering \fmbadge{fmev}{Evitable}\par}
\begin{center}\begin{tikzpicture}\begin{axis}[fmaxis,width=4.4cm,height=3.8cm,xmin=-1,xmax=3,ymin=-1,ymax=4]\draw[fmguide] ({axis cs:1,0}|-{rel axis cs:0,0}) -- ({axis cs:1,0}|-{rel axis cs:0,1});\addplot[fmcurve,domain=-1:3,samples=160]{x/2+1};\node[fmopen] at (axis cs:1,1.5) {};\node[fmdot] at (axis cs:1,3) {};\node[anchor=west,font=\small] at (axis cs:1.05,-0.6) {$a$};\end{axis}\end{tikzpicture}\end{center}

          {\centering\footnotesize Existeix \fmm{\lim\limits_{x\to a}f(x)=k}, però \fmm{\nexists f(a)} o \fmm{f(a)\neq k}.\par}\end{minipage}\hfill\begin{minipage}[t]{0.235\linewidth}\centering {\centering \fmbadge{fmsf}{De salt finit}\par}
\begin{center}\begin{tikzpicture}\begin{axis}[fmaxis,width=4.4cm,height=3.8cm,xmin=-1,xmax=3,ymin=-1,ymax=4]\draw[fmguide] ({axis cs:1,0}|-{rel axis cs:0,0}) -- ({axis cs:1,0}|-{rel axis cs:0,1});\addplot[fmcurve,domain=-1:1,samples=160]{x+1};\addplot[fmcurve,domain=1:3,samples=160]{x/2};\node[fmopen] at (axis cs:1,2) {};\node[fmdot] at (axis cs:1,0.5) {};\node[anchor=west,font=\small] at (axis cs:1.05,-0.6) {$a$};\end{axis}\end{tikzpicture}\end{center}

          {\centering\footnotesize Els dos límits laterals són nombres, però diferents.\par}\end{minipage}\hfill\begin{minipage}[t]{0.235\linewidth}\centering {\centering \fmbadge{fmsi}{De salt infinit}\par}
\begin{center}\begin{tikzpicture}\begin{axis}[fmaxis,width=4.4cm,height=3.8cm,xmin=-1,xmax=3,ymin=-2,ymax=5]\draw[fmasy] ({axis cs:1,0}|-{rel axis cs:0,0}) -- ({axis cs:1,0}|-{rel axis cs:0,1});\addplot[fmcurve,domain=-1:0.9,samples=300]{1/(2*(1-x))};\addplot[fmcurve,domain=1:3,samples=160]{x-2};\node[fmdot] at (axis cs:1,-1) {};\node[anchor=west,font=\small] at (axis cs:1.05,-1.6) {$a$};\end{axis}\end{tikzpicture}\end{center}

          {\centering\footnotesize Un límit lateral és \fmm{\pm\infty} i l'altre és un nombre.\par}\end{minipage}\hfill\begin{minipage}[t]{0.235\linewidth}\centering {\centering \fmbadge{fmas}{Asimptòtica}\par}
\begin{center}\begin{tikzpicture}\begin{axis}[fmaxis,width=4.4cm,height=3.8cm,xmin=-1,xmax=3,ymin=-5,ymax=5]\draw[fmasy] ({axis cs:1,0}|-{rel axis cs:0,0}) -- ({axis cs:1,0}|-{rel axis cs:0,1});\addplot[fmcurve,domain=-1:0.8,samples=300]{1/(x-1)};\addplot[fmcurve,domain=1.2:3,samples=300]{1/(x-1)};\node[anchor=west,font=\small] at (axis cs:1.05,-4.4) {$a$};\end{axis}\end{tikzpicture}\end{center}

          {\centering\footnotesize Els dos límits laterals són \fmm{\pm\infty}.\par}\end{minipage}\par\medskip

\begin{fmnote}Es diu \textbf{evitable} perquè, redefinint \fmm{f(a)=k}, la funció es faria contínua. En les altres (no evitables) el límit no existeix.\end{fmnote}
\end{fmprop}

\Needspace{7\baselineskip}
\begin{fmex}{Exemple: la gràfica de l'apartat 5.5}
\begin{center}\begin{tikzpicture}\begin{axis}[fmaxis,width=8cm,height=6cm,xmin=-4,xmax=6,ymin=-4,ymax=6]\draw[fmasy] ({axis cs:-1,0}|-{rel axis cs:0,0}) -- ({axis cs:-1,0}|-{rel axis cs:0,1});\addplot[fmcurve,domain=-4:-1.1,samples=300]{-1/(x+1)};\addplot[fmcurve,domain=-1:2,samples=160]{x};\addplot[fmcurve,domain=2:6,samples=160]{4-(x-2)^2/2};\node[fmopen] at (axis cs:-1,-1) {};\node[fmopen] at (axis cs:1,1) {};\node[fmdot] at (axis cs:1,3) {};\node[fmopen] at (axis cs:2,2) {};\node[fmdot] at (axis cs:2,4) {};\end{axis}\end{tikzpicture}\end{center}

\begin{itemize}
\item \fmm{x=-1}: \fmm{\nexists f(-1)}, \fmm{\lim_{x\to-1^-}f=+\infty} i \fmm{\lim_{x\to-1^+}f=-1} $\to$ discontinuïtat \fmbadge{fmsi}{de salt infinit}.
\item \fmm{x=1}: \fmm{f(1)=3}, però \fmm{\lim_{x\to1}f=1\neq3} $\to$ discontinuïtat \fmbadge{fmev}{evitable}.
\item \fmm{x=2}: \fmm{\lim_{x\to2^-}f=2} i \fmm{\lim_{x\to2^+}f=4} $\to$ discontinuïtat \fmbadge{fmsf}{de salt finit} (salt de 2 unitats).
\item En la resta de punts, \fmm{f} és contínua.
\end{itemize}
\end{fmex}

\Needspace{18\baselineskip}
\section{Estudi de la continuïtat: funcions racionals i funcions a trossos de tota mena}\hypertarget{sec10}{}

\Needspace{7\baselineskip}
\begin{fmprop}{Continuïtat de les funcions elementals}
\begin{center}\renewcommand{\arraystretch}{1.5}\small\begin{tabularx}{\linewidth}{>{\raggedright\arraybackslash}p{5.2cm}>{\raggedright\arraybackslash}X}\toprule
\textbf{Funció} & \textbf{És contínua\dots{}}\\\midrule
Polinòmiques, $e^x$, $a^x$, $\sin x$, $\cos x$, $\sqrt[3]{x}$, $|x|$ & en tot $\mathbb R$\\[3pt]
$\dfrac{P(x)}{Q(x)}$ & excepte on $Q(x)=0$\\[3pt]
$\sqrt{g(x)}$ (índex parell) & on $g(x)\geq0$\\[3pt]
$\ln g(x)$, $\log_a g(x)$ & on $g(x)\gt0$\\[3pt]
$\tan x=\dfrac{\sin x}{\cos x}$ & excepte en $x=\frac\pi2+k\pi$, on $\cos x=0$ (asímptotes verticals)\\[3pt]
$e^{g(x)}$, $\sin g(x)$, $\cos g(x)$ & on ho siga $g$\\[3pt]
\bottomrule\end{tabularx}\end{center}

Les sumes, restes, productes i composicions de funcions contínues són contínues; els quocients, també, excepte on s'anu\lgem{}a el denominador.
\end{fmprop}

\Needspace{16\baselineskip}
\subsection*{Funcions racionals}

\Needspace{7\baselineskip}
\begin{fmdef}{Mètode}
\begin{enumerate}
\item Busquem els punts que \textbf{anu\lgem{}en el denominador}: \fmm{Q(x)=0}. Fora d'ells, la funció és contínua.
\item En cada punt \fmm{a}, calculem \fmm{\displaystyle\lim_{x\to a}f(x)}: si dona \fmm{\frac00} i en simplificar ix un nombre, la discontinuïtat és \fmbadge{fmev}{evitable}; si dona \fmm{\frac{k}{0}}, és \fmbadge{fmas}{asimptòtica} (asímptota vertical).
\end{enumerate}
\end{fmdef}

\Needspace{7\baselineskip}
\begin{fmex}{Exemple}
Estudiem la continuïtat de \fmm{f(x)=\dfrac{x^2-1}{x^2-3x+2}}.

\fmm{x^2-3x+2=0\Rightarrow x=\dfrac{3\pm\sqrt{9-8}}{2}\Rightarrow x=1,\ x=2}. \fmm{f} és contínua en \fmm{\mathbb R-\{1,2\}}.

\fmm{x=1}: \fmm{\displaystyle\lim_{x\to1}\frac{x^2-1}{x^2-3x+2}\overset{\frac00}{=}\lim_{x\to1}\frac{\cancel{(x-1)}(x+1)}{\cancel{(x-1)}(x-2)}=\frac{2}{-1}=-2}. Com que \fmm{\nexists f(1)}: discontinuïtat \fmbadge{fmev}{evitable}.

\fmm{x=2}: \fmm{\displaystyle\lim_{x\to2}\frac{x+1}{x-2}=\frac30}; laterals: \fmm{-\infty} (amb \fmm{1{,}9}) i \fmm{+\infty} (amb \fmm{2{,}1}): discontinuïtat \fmbadge{fmas}{asimptòtica}.

\begin{center}\begin{tikzpicture}\begin{axis}[fmaxis,width=6.4cm,height=4.6cm,xmin=-3,xmax=5,ymin=-6,ymax=6]\draw[fmasy] ({axis cs:2,0}|-{rel axis cs:0,0}) -- ({axis cs:2,0}|-{rel axis cs:0,1});\draw[fmasy] ({rel axis cs:0,0}|-{axis cs:0,1}) -- ({rel axis cs:1,0}|-{axis cs:0,1});\addplot[fmcurve,domain=-3:1.75,samples=300]{(x+1)/(x-2)};\addplot[fmcurve,domain=2.25:5,samples=300]{(x+1)/(x-2)};\node[fmopen] at (axis cs:1,-2) {};\end{axis}\end{tikzpicture}\end{center}

\begin{fmnote}Si el denominador no s'anu\lgem{}a mai, la funció és contínua en tot \fmm{\mathbb R}: per exemple, \fmm{\dfrac{x+3}{x^2+1}}.\end{fmnote}
\end{fmex}

\Needspace{16\baselineskip}
\subsection*{Funcions a trossos}

\Needspace{7\baselineskip}
\begin{fmdef}{Mètode}
\begin{enumerate}
\item \textbf{Estudiem la continuïtat de cada tros per separat}, en el seu interval obert, amb la taula de dalt. Un problema d'un tros (denominador nul, logaritme d'un nombre \fmm{\leq0}, \fmm{\tan} en \fmm{\frac\pi2}\dots{}) només compta \textbf{si el punt és dins del seu interval}.
\item \textbf{Mirem ara els punts on la funció es trenca} (on canvia l'expressió) i comprovem les tres condicions: \fmm{f(a)}, \fmm{\lim_{x\to a^-}f} (tros de l'esquerra) i \fmm{\lim_{x\to a^+}f} (tros de la dreta).
\item \textbf{Conclusió}: on és contínua i quin tipus de discontinuïtat té en cada punt.
\end{enumerate}

\begin{fmnote}Una funció amb \textbf{valor absolut} es transforma primer en una funció a trossos: \fmm{|x-2|=\begin{cases}-(x-2)&\text{si } x\lt2\\x-2&\text{si } x\geq2\end{cases}}.\end{fmnote}
\end{fmdef}

\Needspace{7\baselineskip}
\begin{fmex}{Exemples resolts}
\begin{enumerate}[label=\textbf{\arabic*.},leftmargin=*]
\item Estudia la continuïtat de
\fmdisp{f(x)=\begin{cases}\dfrac{\sin x}{x}&\text{si } x\lt0\\[4pt]x^2+1&\text{si } 0\leq x\lt2\\[2pt]\sqrt{x+7}+2&\text{si } x\geq2\end{cases}}

Trossos: \fmm{\frac{\sin x}{x}} només falla en \fmm{x=0}, que no és de \fmm{x\lt0}; el segon és polinòmic; \fmm{\sqrt{x+7}} necessita \fmm{x\geq-7}, i en \fmm{x\geq2} es compleix. Queden \fmm{x=0} i \fmm{x=2}.

\fmm{x=0}: \fmm{\ f(0)=1}; \fmm{\ \displaystyle\lim_{x\to0^-}\frac{\sin x}{x}=1}; \fmm{\ \displaystyle\lim_{x\to0^+}(x^2+1)=1} $\to$ \fmbadge{fmok}{contínua}

\fmm{x=2}: \fmm{\ f(2)=\sqrt9+2=5}; \fmm{\ \displaystyle\lim_{x\to2^-}(x^2+1)=5} $\to$ \fmbadge{fmok}{contínua}. \fmm{f} és contínua en tot \fmm{\mathbb R}.
\item Estudia la continuïtat de
\fmdisp{f(x)=\begin{cases}\ln(2-x)&\text{si } x\lt1\\[2pt]\tan\dfrac{\pi x}{4}&\text{si } 1\leq x\lt3\\[4pt]\dfrac{x-4}{x-2}&\text{si } x\geq3\end{cases}}

\fmm{\ln(2-x)} necessita \fmm{x\lt2}: bé en \fmm{x\lt1}. \fmm{\tan\frac{\pi x}{4}} no existeix quan \fmm{\frac{\pi x}{4}=\frac\pi2}, és a dir, en \fmm{x=2}, que \textbf{sí} és de \fmm{[1,3)}. \fmm{\frac{x-4}{x-2}} falla en \fmm{x=2}, que no és de \fmm{x\geq3}.

\fmm{x=1}: \fmm{\ f(1)=\tan\frac\pi4=1}; \fmm{\ \displaystyle\lim_{x\to1^-}\ln(2-x)=\ln1=0} $\to$ \fmbadge{fmsf}{salt finit}

\fmm{x=2}: \fmm{\ \displaystyle\lim_{x\to2^-}\tan\frac{\pi x}{4}=+\infty}, \fmm{\ \displaystyle\lim_{x\to2^+}\tan\frac{\pi x}{4}=-\infty} $\to$ \fmbadge{fmas}{asimptòtica}

\fmm{x=3}: \fmm{\ f(3)=\frac{-1}{1}=-1}; \fmm{\ \displaystyle\lim_{x\to3^-}\tan\frac{3\pi}{4}=-1} $\to$ \fmbadge{fmok}{contínua}. \fmm{f} és contínua en \fmm{\mathbb R-\{1,2\}}.
\item \fmpau{PAU Castella-la Manxa 2020}Determina on és contínua la funció
\fmdisp{f(x)=\begin{cases}\dfrac{3}{x-2}&\text{si } x\lt2\\[4pt]\cos(\pi x)&\text{si } 2\leq x\leq3\\[2pt]\dfrac{\ln(x-2)}{3-x}&\text{si } x\gt3\end{cases}}

Els denominadors s'anu\lgem{}en en \fmm{x=2} i \fmm{x=3}, fora dels seus trossos; \fmm{\ln(x-2)} existeix per a \fmm{x\gt2}. Estudiem \fmm{x=2} i \fmm{x=3}.

\fmm{x=2}: \fmm{\ f(2)=\cos2\pi=1}; \fmm{\ \displaystyle\lim_{x\to2^-}\frac{3}{x-2}=\frac{3}{0^-}=-\infty} $\to$ \fmbadge{fmsi}{salt infinit}

\fmm{x=3}: \fmm{\ f(3)=\cos3\pi=-1}; \fmm{\ \displaystyle\lim_{x\to3^+}\frac{\ln(x-2)}{3-x}\overset{\frac00}{=}} amb \fmm{t=x-3\to0^+}: \fmm{\ \displaystyle\lim_{t\to0^+}\frac{\ln(1+t)}{-t}=-1} $\to$ \fmbadge{fmok}{contínua}

\fmm{f} és contínua en \fmm{\mathbb R-\{2\}}.
\item Amb exponencial, arrel i valor absolut:
\fmdisp{f(x)=\begin{cases}e^{1/x}&\text{si } x\lt0\\\sqrt{x}&\text{si } 0\leq x\leq4\\|x-6|&\text{si } x\gt4\end{cases}}

\fmm{x=0}: \fmm{\ f(0)=0}; \fmm{\ \displaystyle\lim_{x\to0^-}e^{1/x}=e^{-\infty}=0}; \fmm{\ \displaystyle\lim_{x\to0^+}\sqrt x=0} $\to$ \fmbadge{fmok}{contínua} (compte: per la dreta, \fmm{e^{1/x}} aniria a \fmm{+\infty}, però aquest tros només val per a \fmm{x\lt0}).

\fmm{x=4}: \fmm{\ f(4)=2}; \fmm{\ \displaystyle\lim_{x\to4^+}|x-6|=2} $\to$ \fmbadge{fmok}{contínua}. El valor absolut és continu (el «pic» de \fmm{x=6} no és cap salt). \fmm{f} és contínua en \fmm{\mathbb R}.
\end{enumerate}
\end{fmex}

\Needspace{18\baselineskip}
\section{Funcions a trossos que depenen d'un paràmetre}\hypertarget{sec11}{}

\Needspace{7\baselineskip}
\begin{fmdef}{Mètode}
Volem trobar el valor (o valors) del paràmetre que fan la funció contínua.

\begin{enumerate}
\item Comprovem que cada tros és continu en el seu interval (compte amb els denominadors!).
\item En cada punt de canvi \fmm{a}, calculem \fmm{\lim_{x\to a^-}f} i \fmm{\lim_{x\to a^+}f}: dependran del paràmetre.
\item \textbf{Igualem} els límits laterals (i \fmm{f(a)}): obtenim una equació (o un sistema, si hi ha dos paràmetres) i la resolem.
\end{enumerate}
\end{fmdef}

\Needspace{7\baselineskip}
\begin{fmex}{Exemples}
\begin{enumerate}[label=\textbf{\arabic*.},leftmargin=*]
\item Troba \fmm{a} perquè \fmm{f} siga contínua en \fmm{\mathbb R}: \fmm{\ f(x)=\begin{cases}\dfrac{\sin(ax)}{x}&\text{si } x\lt0\\[4pt]x^2+3&\text{si } x\geq0\end{cases}}

Els dos trossos són continus en els seus intervals (\fmm{x=0} no és de \fmm{x\lt0}). \fmm{\ f(0)=\displaystyle\lim_{x\to0^+}(x^2+3)=3}; \fmm{\ \displaystyle\lim_{x\to0^-}\frac{\sin(ax)}{x}=\lim_{x\to0^-}\frac{ax}{x}=a}

\fmm{a=3}
\item Troba \fmm{a} i \fmm{b} perquè \fmm{f} siga contínua en \fmm{\mathbb R}: \fmm{\ f(x)=\begin{cases}-2\sin x&\text{si } x\leq-\frac\pi2\\a\sin x+b&\text{si } -\frac\pi2\lt x\lt\frac\pi2\\\cos x&\text{si } x\geq\frac\pi2\end{cases}}

\fmm{x=-\frac\pi2}: \fmm{\ f\left(-\frac\pi2\right)=-2\cdot(-1)=2}; \fmm{\ \displaystyle\lim_{x\to-\frac\pi2^+}(a\sin x+b)=-a+b\Rightarrow-a+b=2}

\fmm{x=\frac\pi2}: \fmm{\ \displaystyle\lim_{x\to\frac\pi2^-}(a\sin x+b)=a+b}; \fmm{\ f\left(\frac\pi2\right)=\cos\frac\pi2=0\Rightarrow a+b=0}

\fmm{\begin{cases}-a+b=2\\a+b=0\end{cases}\Rightarrow b=1,\ a=-1}
\item Troba \fmm{a} i \fmm{b} perquè \fmm{f} siga contínua en \fmm{\mathbb R}: \fmm{\ f(x)=\begin{cases}a\,e^{x}&\text{si } x\leq0\\[2pt]\dfrac{\ln(1+3x)}{x}&\text{si } 0\lt x\leq1\\[4pt]b\sqrt{x+3}&\text{si } x\gt1\end{cases}}

\fmm{x=0}: \fmm{\ f(0)=a}; \fmm{\ \displaystyle\lim_{x\to0^+}\frac{\ln(1+3x)}{x}=\lim_{x\to0^+}\frac{3x}{x}=3\Rightarrow a=3}

\fmm{x=1}: \fmm{\ f(1)=\ln4}; \fmm{\ \displaystyle\lim_{x\to1^+}b\sqrt{x+3}=2b\Rightarrow b=\frac{\ln4}{2}=\ln2}
\item \fmpau{PAU Aragó 2022}Troba \fmm{a} perquè \fmm{f} siga contínua en \fmm{\mathbb R}: \fmm{\ f(x)=\begin{cases}\sqrt{-x}+e^{3}&\text{si } x\leq0\\(1-x)^{a/x}&\text{si } x\gt0\end{cases}}

\fmm{f(0)=\displaystyle\lim_{x\to0^-}f(x)=e^3}; \fmm{\ \displaystyle\lim_{x\to0^+}(1-x)^{a/x}\overset{1^\infty}{=}e^{\lim\limits_{x\to0^+}(-x)\frac ax}=e^{-a}}

\fmm{e^{-a}=e^{3}\Rightarrow a=-3}
\item \fmpau{PAU Castella i Lleó 2024}Quins valors han de prendre \fmm{a} i \fmm{b\neq0} perquè \fmm{f} siga contínua en \fmm{\mathbb R}? \fmm{\ f(x)=\begin{cases}\dfrac{\cos x-a}{bx^2}&\text{si } x\neq0\\[4pt]1&\text{si } x=0\end{cases}}

El denominador tendeix a 0: perquè el límit siga finit, el numerador també, \fmm{1-a=0\Rightarrow a=1}.

\fmm{\displaystyle\lim_{x\to0}\frac{\cos x-1}{bx^2}=\lim_{x\to0}\frac{-\frac{x^2}{2}}{bx^2}=-\frac{1}{2b}=1\Rightarrow b=-\frac12}
\end{enumerate}
\end{fmex}

\Needspace{18\baselineskip}
\section{Teorema de Bolzano i propietats de les funcions contínues}\hypertarget{sec12}{}

\Needspace{7\baselineskip}
\begin{fmprop}{Teorema de Bolzano}
Si \fmm{f} és \textbf{contínua en l'interval tancat \fmm{[a,b]}} i \fmm{f(a)} i \fmm{f(b)} tenen \textbf{signes contraris} (\fmm{f(a)\cdot f(b)\lt0}), aleshores existeix almenys un \fmm{c\in(a,b)} tal que

\fmdisp{f(c)=0}

\textbf{Interpretació geomètrica}: si una corba contínua comença per davall de l'eix \fmm{OX} i acaba per damunt (o al revés), l'ha de tallar almenys una vegada. Una funció contínua no pot canviar de signe «botant» el zero.
\end{fmprop}

\Needspace{7\baselineskip}
\begin{fmdef}{Què diu i què no diu}
\begin{itemize}
\item Assegura que hi ha \textbf{almenys una} arrel, però no diu quantes n'hi ha ni on és exactament. (Per veure que és \textbf{única}, cal saber que \fmm{f} és estrictament creixent o decreixent en l'interval: ho farem amb derivades a la unitat següent.)
\item \textbf{La continuïtat és imprescindible}: \fmm{f(x)=\frac1x} compleix \fmm{f(-1)=-1\lt0} i \fmm{f(1)=1\gt0}, però no s'anu\lgem{}a mai: no és contínua en \fmm{[-1,1]}.
\item El recíproc és fals: si \fmm{f(a)} i \fmm{f(b)} tenen el mateix signe, la funció pot tindre arrels o no (\fmm{x^2-1} en \fmm{[-2,2]} en té dues).
\end{itemize}
\end{fmdef}

\Needspace{7\baselineskip}
\begin{fmprop}{Conseqüències}
\begin{itemize}
\item \textbf{Teorema dels valors intermedis (Darboux)}: si \fmm{f} és contínua en \fmm{[a,b]}, pren \textbf{tots} els valors compresos entre \fmm{f(a)} i \fmm{f(b)}. Prova: si \fmm{k} està entre \fmm{f(a)} i \fmm{f(b)}, apliquem Bolzano a \fmm{g(x)=f(x)-k}.
\item \textbf{Dues corbes es tallen}: per veure que \fmm{f(x)=g(x)} té solució en \fmm{(a,b)}, apliquem Bolzano a \fmm{h(x)=f(x)-g(x)}.
\item \textbf{Teorema de Weierstrass}: si \fmm{f} és contínua en un interval \textbf{tancat} \fmm{[a,b]}, hi assoleix un \textbf{màxim absolut} i un \textbf{mínim absolut}. (En un interval obert pot fallar: \fmm{\frac1x} en \fmm{(0,1)} no té màxim.)
\end{itemize}
\end{fmprop}

\Needspace{7\baselineskip}
\begin{fmex}{Exemples}
\begin{enumerate}[label=\textbf{\arabic*.},leftmargin=*]
\item Demostra que l'equació \fmm{x^3+x-1=0} té almenys una solució en \fmm{(0,1)} i troba-la amb una xifra decimal.

\fmm{f(x)=x^3+x-1} és polinòmica: contínua en \fmm{[0,1]}. \fmm{f(0)=-1\lt0}, \fmm{f(1)=1\gt0}. Per Bolzano, \fmm{\exists c\in(0,1)} amb \fmm{f(c)=0}.

\textbf{Bisecció}: \fmm{f(0{,}5)=-0{,}375\lt0\Rightarrow c\in(0{,}5;\,1)}; \fmm{\ f(0{,}7)=0{,}043\gt0\Rightarrow c\in(0{,}5;\,0{,}7)}; \fmm{\ f(0{,}6)=-0{,}184\lt0\Rightarrow c\in(0{,}6;\,0{,}7)}. Per tant \fmm{c\approx0{,}6\ldots} (de fet, \fmm{c\approx0{,}682}).
\item Demostra que les gràfiques de \fmm{y=\cos x} i \fmm{y=x} es tallen en algun punt de \fmm{\left(0,\frac\pi2\right)}.

Sigui \fmm{h(x)=\cos x-x}, contínua en \fmm{\mathbb R}. \fmm{h(0)=1\gt0}, \fmm{h\left(\frac\pi2\right)=0-\frac\pi2\lt0}.

Per Bolzano, \fmm{\exists c\in\left(0,\frac\pi2\right)} amb \fmm{h(c)=0}, és a dir, \fmm{\cos c=c}: les corbes es tallen en \fmm{x=c}.
\item Demostra que \fmm{f(x)=x^3-x} pren el valor \fmm{5} en algun punt de \fmm{[1,2]}.

\fmm{f} és contínua, \fmm{f(1)=0} i \fmm{f(2)=6}. Com que \fmm{0\lt5\lt6}, pel teorema dels valors intermedis \fmm{\exists c\in(1,2)} amb \fmm{f(c)=5}. (O bé: Bolzano a \fmm{g(x)=x^3-x-5}, amb \fmm{g(1)=-5}, \fmm{g(2)=1}.)
\item \fmpau{PAU Catalunya 2019}Sigui \fmm{f(x)=\dfrac{2x^3-5x+4}{1-x}}. Sabem que \fmm{f(-2)=-\frac23}, \fmm{f(0)=4} i \fmm{f(2)=-10}. Podem deduir que \fmm{(-2,0)} conté un zero de \fmm{f}? I \fmm{(0,2)}?

\fmm{f} és contínua en \fmm{\mathbb R-\{1\}} (en \fmm{x=1} ix \fmm{\frac{1}{0}}: asímptota vertical).

En \fmm{[-2,0]} és contínua i canvia de signe: \textbf{sí}, hi ha un zero. En \fmm{[0,2]} també canvia de signe, però \textbf{no} és contínua (falla en \fmm{x=1}): Bolzano no es pot aplicar i no podem deduir res. (De fet, el canvi de signe el provoca l'asímptota.)

Afinant: \fmm{f(-1)=\frac{-2+5+4}{2}=\frac72\gt0} i \fmm{f(-2)\lt0}: el zero és en \fmm{(-2,-1)}.
\end{enumerate}
\end{fmex}

\Needspace{18\baselineskip}
\section{Asímptotes}\hypertarget{sec13}{}

\Needspace{7\baselineskip}
\begin{fmdef}{Asímptota}
Una \textbf{asímptota} és una recta a la qual la gràfica d'una funció s'acosta cada vegada més, sense arribar a tocar-la, quan ens n'allunyem cap a l'infinit. Les asímptotes es calculen amb \textbf{límits}: són la «traducció gràfica» dels límits infinits i dels límits a l'infinit que hem estudiat en aquesta unitat.

\begin{itemize}
\item \textbf{Verticals}: rectes \fmm{x=a}.
\item \textbf{Horitzontals}: rectes \fmm{y=b}.
\item \textbf{Obliqües}: rectes \fmm{y=mx+n}, amb \fmm{m\neq0}.
\end{itemize}
\end{fmdef}

\Needspace{16\baselineskip}
\subsection*{Asímptotes verticals}

\Needspace{7\baselineskip}
\begin{fmdef}{Asímptota vertical}
\par\noindent \begin{minipage}[t]{0.48\linewidth}\vspace{0pt}\raggedright La recta \fmm{x=a} és una \textbf{asímptota vertical} (AV) de \fmm{f} si almenys un dels límits laterals en \fmm{a} és infinit:

\fmdisp{\lim_{x\to a^-}f(x)=\pm\infty\quad o\quad\lim_{x\to a^+}f(x)=\pm\infty}

És el que en l'apartat 5.10 hem anomenat una discontinuïtat \fmbadge{fmas}{asimptòtica}.\end{minipage}\hfill\begin{minipage}[t]{0.48\linewidth}\vspace{0pt}\raggedright \begin{center}\begin{tikzpicture}\begin{axis}[fmaxis,width=6.4cm,height=4.6cm,xmin=-5,xmax=5,ymin=-6,ymax=6]\draw[fmasy] ({axis cs:-2,0}|-{rel axis cs:0,0}) -- ({axis cs:-2,0}|-{rel axis cs:0,1});\draw[fmasy] ({axis cs:2,0}|-{rel axis cs:0,0}) -- ({axis cs:2,0}|-{rel axis cs:0,1});\draw[fmasy] ({rel axis cs:0,0}|-{axis cs:0,0}) -- ({rel axis cs:1,0}|-{axis cs:0,0});\addplot[fmcurve,domain=-5:-2.02,samples=600]{(x+1)/(x^2-4)};\addplot[fmcurve,domain=-1.985:1.985,samples=900]{(x+1)/(x^2-4)};\addplot[fmcurve,domain=2.02:5,samples=600]{(x+1)/(x^2-4)};\end{axis}\end{tikzpicture}\end{center}\end{minipage}\par\medskip

\fmuline{On les busquem?} En una funció racional \fmm{\dfrac{P(x)}{Q(x)}}, en els punts on \fmm{Q(x)=0}:

\begin{itemize}
\item Si en substituir ix \fmm{\dfrac{k}{0}}, amb \fmm{k\neq0}: \textbf{hi ha AV}. Els signes dels límits laterals ens diuen si la gràfica puja (\fmm{+\infty}) o baixa (\fmm{-\infty}) a cada costat.
\item Si ix \fmm{\dfrac00}: simplifiquem. Si el límit dona un nombre, \textbf{no hi ha AV} (és una discontinuïtat evitable, un «forat»); si torna a donar \fmm{\dfrac k0}, sí que n'hi ha.
\end{itemize}

\begin{fmnote}Les funcions polinòmiques i les exponencials del tipus \fmm{a^x} no tenen asímptotes verticals. Una gràfica no pot tallar mai una asímptota vertical.\end{fmnote}
\end{fmdef}

\Needspace{7\baselineskip}
\begin{fmex}{Exemples}
\begin{enumerate}[label=\textbf{\arabic*.},leftmargin=*]
\item Asímptotes verticals de \fmm{f(x)=\dfrac{x+1}{x^2-4}} (és la funció del gràfic).

\fmm{x^2-4=0\Rightarrow x=-2,\ x=2}. En tots dos el numerador no s'anu\lgem{}a: \fmm{\frac{-1}{0}} i \fmm{\frac{3}{0}}.

\fmm{x=-2}: \fmm{\ \displaystyle\lim_{x\to-2^-}\frac{x+1}{x^2-4}=\frac{-1}{0^+}=-\infty}, \fmm{\ \displaystyle\lim_{x\to-2^+}\frac{x+1}{x^2-4}=\frac{-1}{0^-}=+\infty} (amb \fmm{-2{,}1} i \fmm{-1{,}9}).

\fmm{x=2}: \fmm{\ \displaystyle\lim_{x\to2^-}\frac{x+1}{x^2-4}=\frac{3}{0^-}=-\infty}, \fmm{\ \displaystyle\lim_{x\to2^+}\frac{x+1}{x^2-4}=\frac{3}{0^+}=+\infty}.

\textbf{AV: \fmm{x=-2} i \fmm{x=2}}.
\item Asímptotes verticals de \fmm{f(x)=\dfrac{x^2-1}{x^2-3x+2}} (l'exemple de l'apartat 5.10).

El denominador s'anu\lgem{}a en \fmm{x=1} i en \fmm{x=2}.

\fmm{x=1}: \fmm{\frac00\Rightarrow\displaystyle\lim_{x\to1}\frac{(x-1)(x+1)}{(x-1)(x-2)}=-2}. Límit finit: \textbf{no} hi ha AV (forat en \fmm{(1,-2)}).

\fmm{x=2}: \fmm{\frac30\Rightarrow} \textbf{AV \fmm{x=2}}.
\end{enumerate}
\end{fmex}

\Needspace{16\baselineskip}
\subsection*{Asímptotes horitzontals}

\Needspace{7\baselineskip}
\begin{fmdef}{Asímptota horitzontal}
\par\noindent \begin{minipage}[t]{0.48\linewidth}\vspace{0pt}\raggedright La recta \fmm{y=b} és una \textbf{asímptota horitzontal} (AH) de \fmm{f} si el límit a l'infinit és un nombre:

\fmdisp{\lim_{x\to+\infty}f(x)=b\quad o\quad\lim_{x\to-\infty}f(x)=b}

Cal calcular els dos límits (\fmm{+\infty} i \fmm{-\infty}): una funció pot tindre AH només per un costat, o una AH diferent per cada costat.\end{minipage}\hfill\begin{minipage}[t]{0.48\linewidth}\vspace{0pt}\raggedright \begin{center}\begin{tikzpicture}\begin{axis}[fmaxis,width=6.4cm,height=4.6cm,xmin=-8,xmax=8,ymin=-1.5,ymax=4.5]\draw[fmasy] ({rel axis cs:0,0}|-{axis cs:0,3}) -- ({rel axis cs:1,0}|-{axis cs:0,3});\addplot[fmcurve,domain=-8:8,samples=200]{(3*x^2-1)/(x^2+2)};\end{axis}\end{tikzpicture}\end{center}\end{minipage}\par\medskip

\fmuline{Funcions racionals}: comparem els graus, com en l'apartat 5.2.

\begin{itemize}
\item grau \fmm{P\lt} grau \fmm{Q\ \Rightarrow} AH \fmm{y=0}.
\item grau \fmm{P=} grau \fmm{Q\ \Rightarrow} AH \fmm{y=\dfrac{a_n}{b_n}} (quocient dels coeficients principals).
\item grau \fmm{P\gt} grau \fmm{Q\ \Rightarrow} no té AH (el límit és infinit).
\end{itemize}

\begin{fmnote}\fmuline{Posició}: el signe de \fmm{f(x)-b} per a \fmm{x} molt gran (o molt negatiu) ens diu si la gràfica queda per damunt (\fmm{+}) o per davall (\fmm{-}) de l'asímptota. A diferència de les verticals, la gràfica sí que pot tallar una asímptota horitzontal.\end{fmnote}
\end{fmdef}

\Needspace{7\baselineskip}
\begin{fmex}{Exemples}
\begin{enumerate}[label=\textbf{\arabic*.},leftmargin=*]
\item \fmm{f(x)=\dfrac{3x^2-1}{x^2+2}} (és la funció del gràfic).

Graus iguals: \fmm{\displaystyle\lim_{x\to\pm\infty}\frac{3x^2-1}{x^2+2}=\frac31=3}. \textbf{AH: \fmm{y=3}} pels dos costats. No té AV: \fmm{x^2+2\neq0}.

Posició: \fmm{f(x)-3=\dfrac{3x^2-1-3x^2-6}{x^2+2}=\dfrac{-7}{x^2+2}\lt0}: la gràfica queda sempre \textbf{per davall} de l'asímptota.
\item \fmm{f(x)=\dfrac{2x}{x^2+1}}

Grau de baix major: \fmm{\displaystyle\lim_{x\to\pm\infty}f(x)=0\Rightarrow} \textbf{AH: \fmm{y=0}}. Per la dreta \fmm{f(x)\gt0} (per damunt); per l'esquerra \fmm{f(x)\lt0} (per davall).
\item \fmm{f(x)=2^x-1}

\fmm{\displaystyle\lim_{x\to-\infty}(2^x-1)=0-1=-1\Rightarrow} \textbf{AH \fmm{y=-1}} per l'esquerra. \fmm{\ \displaystyle\lim_{x\to+\infty}(2^x-1)=+\infty}: per la dreta no n'hi ha.
\end{enumerate}
\end{fmex}

\Needspace{16\baselineskip}
\subsection*{Asímptotes obliqües: mètode del quocient}

\Needspace{7\baselineskip}
\begin{fmdef}{Asímptota obliqua}
\par\noindent \begin{minipage}[t]{0.48\linewidth}\vspace{0pt}\raggedright La recta \fmm{y=mx+n}, amb \fmm{m\neq0}, és una \textbf{asímptota obliqua} (AO) si la gràfica s'hi acosta quan \fmm{x\to\pm\infty}:

\fmdisp{\lim_{x\to\pm\infty}\big[f(x)-(mx+n)\big]=0}

Una funció racional \fmm{\dfrac{P(x)}{Q(x)}} té asímptota obliqua quan \textbf{grau \fmm{P=} grau \fmm{Q+1}} (el numerador té exactament un grau més). En aquest cas no té asímptota horitzontal.\end{minipage}\hfill\begin{minipage}[t]{0.48\linewidth}\vspace{0pt}\raggedright \begin{center}\begin{tikzpicture}\begin{axis}[fmaxis,width=6.4cm,height=4.6cm,xmin=-6,xmax=8,ymin=-8,ymax=12]\draw[fmasy] ({axis cs:1,0}|-{rel axis cs:0,0}) -- ({axis cs:1,0}|-{rel axis cs:0,1});\addplot[fmasy,domain=-6:8]{1*x+1};\addplot[fmcurve,domain=-6:0.7,samples=300]{(x^2+3)/(x-1)};\addplot[fmcurve,domain=1.3:8,samples=300]{(x^2+3)/(x-1)};\end{axis}\end{tikzpicture}\end{center}\end{minipage}\par\medskip

\fmuline{Mètode del quocient}: dividim \fmm{P(x):Q(x)} i obtenim el quocient \fmm{C(x)=mx+n} i el residu \fmm{R(x)}:

\fmdisp{P(x)=Q(x)\cdot C(x)+R(x)\quad\Longrightarrow\quad f(x)=C(x)+\frac{R(x)}{Q(x)}}

Com que \fmm{\dfrac{R(x)}{Q(x)}\to0} quan \fmm{x\to\pm\infty} (grau de dalt menor), la gràfica s'acosta a la recta: \textbf{l'asímptota obliqua és \fmm{y=C(x)}}. El signe de \fmm{\dfrac{R(x)}{Q(x)}} ens diu si la gràfica queda per damunt o per davall.
\end{fmdef}

\Needspace{7\baselineskip}
\begin{fmex}{Exemples}
\begin{enumerate}[label=\textbf{\arabic*.},leftmargin=*]
\item Totes les asímptotes de \fmm{f(x)=\dfrac{x^2+3}{x-1}} (és la funció del gràfic).

AV: \fmm{x-1=0\Rightarrow x=1}, \fmm{\ \frac40}: \fmm{\displaystyle\lim_{x\to1^-}f(x)=-\infty}, \fmm{\ \displaystyle\lim_{x\to1^+}f(x)=+\infty}. \textbf{AV \fmm{x=1}}.

Graus 2 i 1: no té AH (el límit és \fmm{\pm\infty}), però sí AO. Dividim \fmm{(x^2+3):(x-1)} (per Ruffini o amb la divisió llarga):

\begin{center}\begin{tabular}{r|rrr}
 & $1$ & $0$ & $3$\\
$1$ &  & $1$ & $1$\\\hline
 & $1$ & $1$ & \fbox{$4$}\\
\end{tabular}\end{center}

\fmm{C(x)=x+1,\ R=4\Rightarrow f(x)=x+1+\dfrac{4}{x-1}}. \textbf{AO: \fmm{y=x+1}}.

Posició: si \fmm{x\to+\infty}, \fmm{\frac{4}{x-1}\gt0}: per damunt; si \fmm{x\to-\infty}, \fmm{\frac{4}{x-1}\lt0}: per davall.
\item \fmm{f(x)=\dfrac{2x^2-x+1}{x}}

AV \fmm{x=0} (\fmm{\frac10}). La divisió entre un monomi és directa: \fmm{f(x)=2x-1+\dfrac1x\Rightarrow} \textbf{AO \fmm{y=2x-1}}.
\item Té asímptota obliqua \fmm{f(x)=\dfrac{x^3+1}{x}}?

Graus 3 i 1: el numerador té \textbf{dos} graus més. El quocient és \fmm{x^2}, que no és una recta: \textbf{no} té AO (ni AH).
\end{enumerate}
\end{fmex}

\Needspace{7\baselineskip}
\begin{fmprop}{Resum: asímptotes d'una funció racional $\frac{P(x)}{Q(x)}$}
\begin{center}\renewcommand{\arraystretch}{1.5}\small\begin{tabularx}{\linewidth}{>{\raggedright\arraybackslash}p{5.2cm}>{\raggedright\arraybackslash}X>{\raggedright\arraybackslash}X}\toprule
\textbf{Graus} & \textbf{Horitzontal} & \textbf{Obliqua}\\\midrule
grau $P\lt$ grau $Q$ & $y=0$ & no\\[3pt]
grau $P=$ grau $Q$ & $y=\dfrac{a_n}{b_n}$ & no\\[3pt]
grau $P=$ grau $Q+1$ & no & $y=C(x)$, quocient de $P:Q$\\[3pt]
grau $P\geq$ grau $Q+2$ & no & no\\[3pt]
\bottomrule\end{tabularx}\end{center}

\begin{fmnote}Les \textbf{verticals} no depenen dels graus: són en els zeros de \fmm{Q} que donen \fmm{\frac k0}. Per aquest motiu, una funció racional pot tindre moltes AV, però com a molt una AH o una AO.\end{fmnote}
\end{fmprop}

\Needspace{16\baselineskip}
\subsection*{Asímptotes de funcions no racionals}

\Needspace{7\baselineskip}
\begin{fmdef}{Mètode general per a l'asímptota obliqua}
Si la funció no és racional no podem dividir polinomis. La recta \fmm{y=mx+n} (\fmm{m\neq0}) és asímptota obliqua quan \fmm{x\to+\infty} (o \fmm{-\infty}) si existeixen i són finits

\fmdisp{m=\lim_{x\to\pm\infty}\frac{f(x)}{x}\qquad\qquad n=\lim_{x\to\pm\infty}\big(f(x)-mx\big)}

\begin{itemize}
\item Cal estudiar \textbf{cada costat per separat}: amb arrels i exponencials, l'asímptota pot ser diferent a \fmm{+\infty} i a \fmm{-\infty}, o existir només en un dels dos.
\item Les verticals estan on la funció no està definida o als extrems del domini: zeros d'un denominador, punts on l'argument d'un logaritme tendeix a 0, valors on \fmm{\tan} no existeix\dots{} Sempre cal comprovar que el límit és infinit.
\end{itemize}
\end{fmdef}

\Needspace{7\baselineskip}
\begin{fmex}{Exemples}
\begin{enumerate}[label=\textbf{\arabic*.},leftmargin=*]
\item \fmm{f(x)=\sqrt{x^2+2x}} (domini \fmm{(-\infty,-2]\cup[0,+\infty)})

\fmm{x\to+\infty}: \fmm{\ m=\displaystyle\lim\frac{\sqrt{x^2+2x}}{x}=1}, \fmm{\ n=\displaystyle\lim\left(\sqrt{x^2+2x}-x\right)=\lim\frac{2x}{\sqrt{x^2+2x}+x}=1}. AO \fmm{y=x+1}.

\fmm{x\to-\infty}: \fmm{\sqrt{x^2+2x}\approx|x|=-x\Rightarrow m=-1}; \fmm{\ n=\displaystyle\lim_{x\to-\infty}\left(\sqrt{x^2+2x}+x\right)\overset{x=-t}{=}\lim_{t\to+\infty}\frac{-2t}{\sqrt{t^2-2t}+t}=-1}. AO \fmm{y=-x-1}.

\begin{center}\begin{tikzpicture}\begin{axis}[fmaxis,width=6.4cm,height=4.6cm,xmin=-7,xmax=6,ymin=-1,ymax=7]\addplot[fmasy,domain=-7:6]{1*x+1};\addplot[fmasy,domain=-7:6]{-1*x+-1};\addplot[fmcurve,domain=-7:-2,samples=400]{sqrt(max(x^2+2*x,0))};\addplot[fmcurve,domain=0:6,samples=400]{sqrt(max(x^2+2*x,0))};\node[fmdot] at (axis cs:-2,0) {};\node[fmdot] at (axis cs:0,0) {};\end{axis}\end{tikzpicture}\end{center}
\item \fmm{f(x)=e^{1/x}}

\fmm{\displaystyle\lim_{x\to0^+}e^{1/x}=e^{+\infty}=+\infty}, \fmm{\ \displaystyle\lim_{x\to0^-}e^{1/x}=e^{-\infty}=0}: AV \fmm{x=0} només per la dreta.

\fmm{\displaystyle\lim_{x\to\pm\infty}e^{1/x}=e^0=1}: AH \fmm{y=1} pels dos costats.

\begin{center}\begin{tikzpicture}\begin{axis}[fmaxis,width=6.4cm,height=4.6cm,xmin=-6,xmax=6,ymin=-1,ymax=6]\draw[fmasy] ({axis cs:0,0}|-{rel axis cs:0,0}) -- ({axis cs:0,0}|-{rel axis cs:0,1});\draw[fmasy] ({rel axis cs:0,0}|-{axis cs:0,1}) -- ({rel axis cs:1,0}|-{axis cs:0,1});\addplot[fmcurve,domain=-6:-0.05,samples=300]{exp(1/x)};\addplot[fmcurve,domain=0.17:6,samples=400]{exp(1/x)};\node[fmopen] at (axis cs:0,0) {};\end{axis}\end{tikzpicture}\end{center}
\item \fmm{f(x)=\ln(x^2-4)} (domini \fmm{|x|\gt2}): \fmm{\displaystyle\lim_{x\to2^+}\ln(x^2-4)=\ln0^+=-\infty} i igual en \fmm{x=-2^-}: AV \fmm{x=\pm2}. Quan \fmm{x\to\pm\infty}, \fmm{f\to+\infty} i \fmm{\frac{f(x)}{x}\to0}: ni AH ni AO.
\item \fmm{f(x)=x+\sin x}: \fmm{\ m=\displaystyle\lim_{x\to+\infty}\left(1+\frac{\sin x}{x}\right)=1}, però \fmm{n=\displaystyle\lim_{x\to+\infty}\sin x} no existeix: \textbf{no} té asímptota obliqua (la gràfica osci\lgem{}a al voltant de \fmm{y=x}).
\end{enumerate}
\end{fmex}

  
\end{document}
